% =====================================================================
%  Adaptive Hierarchical KV State Management
%  Cost-Asymmetric Hierarchical KV State Management for Long-Context LLM Inference
%
%  編譯方式（必須用 XeLaTeX，因為含中文）：
%     xelatex main
%     bibtex  main
%     xelatex main
%     xelatex main
%
%  ⚠️ 本稿為初稿（draft）。Evaluation 章為「實驗計畫」，
%     所有結果表格均為佔位符，尚未執行實驗。
% =====================================================================
\documentclass[10pt,twocolumn]{article}

% ---------- 中文支援 ----------
% 西文字型（XeLaTeX 下用 fontspec，不要用 times/helvet 套件）
\usepackage{fontspec}
\setmainfont{TeX Gyre Termes}
\setsansfont{TeX Gyre Heros}
\setmonofont{TeX Gyre Cursor}[Scale=0.9]

\usepackage{xeCJK}
% Linux（本機）字型。macOS 上可改回 Songti TC / Heiti TC。
\setCJKmainfont{Noto Serif CJK TC}[AutoFakeSlant=0.15]
\setCJKsansfont{Noto Sans CJK TC}[AutoFakeSlant=0.15]
\setCJKmonofont{Noto Sans Mono CJK TC}
\xeCJKsetup{CJKmath=true}

% ---------- 版面 ----------
\usepackage[letterpaper,margin=0.75in]{geometry}
\usepackage{amsmath,amssymb}
\usepackage{booktabs}
\usepackage{multirow}
\usepackage{array}
\usepackage{graphicx}
\usepackage{xcolor}
\usepackage{tikz}
\usetikzlibrary{arrows.meta,positioning,fit,backgrounds,shapes.geometric,calc,decorations.pathreplacing}
\usepackage[hidelinks]{hyperref}
\usepackage{cite}
\usepackage{caption}
\captionsetup{font=small,labelfont=bf}

% ---------- 自訂 ----------
\newcommand{\sys}{\mbox{Tiara}\xspace} % Tiered, quality-Aware, Recompute-Aware KV manager
\usepackage{xspace}
\definecolor{gpucol}{HTML}{C0392B}
\definecolor{cpucol}{HTML}{2980B9}
\definecolor{ssdcol}{HTML}{27AE60}
\definecolor{dropcol}{HTML}{7F8C8D}
\definecolor{hlbox}{HTML}{FDF2E9}

\newcommand{\todo}[1]{{\color{red}\textbf{[TODO: #1]}}}
\newcommand{\placeholder}[1]{{\color{gray}\textit{#1}}}
\newcommand{\act}[1]{\textsf{\footnotesize #1}}  % 動作名稱

\title{\vspace{-1.5em}\textbf{成本不對稱下以品質為約束的\\
可恢復階層式 KV State 管理}\\[0.3em]
\large Cost-Asymmetric Hierarchical KV State Management\\
for Long-Context LLM Inference}

\author{\textit{（作者資訊於投稿時填入）}}
\date{}

\begin{document}
\maketitle

% =====================================================================
\begin{abstract}
\noindent
長上下文（long-context）LLM 推論的核心系統瓶頸之一，是 KV cache 隨上下文長度線性成長並耗盡 GPU 高頻寬記憶體（HBM）。既有工作已分別處理了這個問題的各個面向：可恢復的驅逐、學習式的未來效用預測、重算與傳輸的取捨，以及以品質為約束的壓縮。然而，據我們所知，尚無工作將這些決策整合到\emph{單一的線上決策}之中。

本文的出發點是一個橫跨硬體的量化事實：\emph{重算與傳輸的成本比 $\kappa$ 隨平台變動達 32 倍}——於 AMD Instinct MI300X（192\,GB HBM3）為 6--21 倍，於 NVIDIA RTX 3090（24\,GB GDDR6X）則為 54--190 倍。此差異使「單一套 KV 放置策略適用於所有硬體」的隱含前提不成立。

問題的\emph{性質}同樣隨平台改變。在 MI300X 上，8B 至 14B 級模型於 128K 上下文的單請求 KV footprint 僅佔 HBM 的 8--10\%，Llama-3.1-8B 甚至可在單卡容納 1M token 的上下文（需 YaRN 外推）；此處的壓力來自多租戶下的\emph{機會成本}，而非容量本身。在 24\,GB 級的加速器上，單請求上下文本身即構成容量壓力，既有文獻的經典設定依然成立。兩者皆為真實的部署情境，而最佳策略在二者上並不相同。

基於此，我們將 KV 放置形式化為一個受品質約束的序列決策問題，其動作空間包含 GPU 全精度、GPU FP8、GPU INT4、CPU、SSD 與「丟棄後重算」六種狀態，並在「品質退化不超過 $\epsilon$」的硬性約束下，聯合最佳化未來效用、傳輸成本、重算成本與記憶體壓力。我們特別論證：重算的價值不在於單次事件的成本（其代價比傳輸高 6--21 倍於 MI300X，54--190 倍於 RTX 3090），而在於它在被需要之前不佔用任何資源；其正確性因而完全取決於未來效用預測的品質。

我們設計了雙平台實驗（RTX 3090 的 30:1 與 MI300X 的 84:1，相距 2.8 倍），使硬體條件化的主張成為可證偽的。二者\emph{並非}此光譜的兩端——GH200 的 9:1 更低，MI355X 的 127:1 更高——但構成我們實際可取得的最大跨度。\textbf{本文為初稿。}第 \ref{sec:eval} 節給出完整的實驗計畫與消融設計，實驗結果尚未取得。
\end{abstract}

\vspace{0.5em}
\noindent\textbf{Keywords:} LLM inference; KV cache; memory hierarchy; long context; quality-constrained optimization; cost-sensitive learning; AMD MI300X; NVIDIA RTX 3090; vLLM

% =====================================================================
\section{Introduction}
\label{sec:intro}

Transformer 模型在自回歸生成時，會將每個已處理 token 的 key 與 value 向量保存下來，避免在生成每個新 token 時重算整段前綴。這組中間狀態即 KV cache，是以記憶體換取計算時間的經典設計。其代價是記憶體佔用隨上下文長度線性成長：對 Llama-3.1-8B（32 層、8 個 KV head、head dimension 128、BF16），每個 token 需要 128\,KiB，1M token 的上下文即需要 122\,GB。在 80\,GB 級的加速器上，這使得系統的限制因素從計算能力轉移到記憶體容量。

圍繞此問題的研究已密集展開。既有工作大致沿四條軸線展開：（i）\textbf{分層儲存}，將 KV 卸載到 CPU DRAM 或 SSD~\cite{flexgen2023,infinigen2024,cachedattention2024,kvdrive2026,tutti2026}；（ii）\textbf{可恢復驅逐}，使被降級的 KV 在重新變得重要時能夠召回~\cite{arkvale2024,qevict2026}；（iii）\textbf{學習式未來效用預測}，以訓練得到的策略取代手寫啟發式規則~\cite{kvp2026,foresightkv2026,lookaheadkv2026}；（iv）\textbf{重算與傳輸的取捨}~\cite{cake2025,kvpr2025,hcache2025}。此外，KVServe~\cite{kvserve2026} 首次將模型品質寫為壓縮決策的硬性約束，而 Fang 等人~\cite{hetmem2025} 首次將分層 KV 放置形式化並推導理論上界。

\paragraph{本文的出發點：一個不再成立的前提。}
上述工作幾乎全部在 NVIDIA A100/H100 級硬體（40--80\,GB HBM）上評估。我們的量化分析（第 \ref{sec:motivation} 節）顯示，在 AMD Instinct MI300X（192\,GB HBM3）上，8B--14B 級模型於 128K 上下文下的單請求 KV footprint 僅 15.6--19.5\,GB，佔 HBM 不到 10\%。\emph{在此類硬體上，既有文獻用以立論的「單請求容量壓力」場景幾乎不存在}。

這並不代表問題消失，而是問題的\emph{性質改變}了。在 192\,GB 的 HBM 上，一個 session 的冷 KV 所佔據的每 GB，都是其他請求無法用來擴大批次的容量。問題因此從「放不下，該驅逐什麼」轉變為「放得下，但該保留嗎」，成為一個機會成本問題，而機會成本只在多租戶場景中存在。

\paragraph{第二個觀察：卸載在大容量 HBM 硬體上更昂貴。}
MI300X 的 HBM 頻寬（5.3\,TB/s 峰值）與主機鏈路頻寬（PCIe Gen5 $\times$16，單向理論 63.0\,GB/s）之比為 84:1，而 H100 為 52:1，MI355X 更達 127:1。此比值持續惡化，因為 HBM 容量與頻寬逐代成長，而主機鏈路自 MI300X 至 MI355X 完全未變。灌滿或抽乾 MI300X 上 Llama-3.1-8B 的 157.8\,GiB KV 預算，經 PCIe 需 2.69 秒，而在 HBM 內僅需 32 毫秒。此比值必然等於頻寬比 84 倍。

\paragraph{貢獻。}
\begin{enumerate}\itemsep0.2em
  \item 我們量化了\emph{重算對傳輸的成本比 $\kappa$ 隨平台變動達 32 倍}（MI300X 上 6--21 倍、RTX 3090 上 54--190 倍），並指出正確的判準是 FLOPS 對主機鏈路之比而非 HBM 頻寬比。同一組量化分析亦顯示，既有 KV 卸載文獻的容量前提在大容量 HBM 加速器上不成立，問題性質由容量轉為機會成本，而在 24\,GB 級加速器上則仍然成立（第 \ref{sec:motivation} 節）。
  \item 我們系統性地比對了 15 篇近期工作，指出四項成本（未來效用、傳輸成本、重算成本、記憶體壓力）從未在單一線上決策中同時出現；現有最高覆蓋為 Marconi~\cite{marconi2025} 的三項，其因僅在 GPU 上運作而結構性地缺少傳輸成本（第 \ref{sec:related} 節）。
  \item 我們將 KV 放置形式化為受品質約束的序列決策問題，動作空間包含六種狀態且明確納入「丟棄後重算」；並論證重算的價值來自其資源佔用特性而非單次成本，因而其正確性取決於未來效用預測（第 \ref{sec:formulation} 節）。
  \item 我們提出 \sys 的設計，透過 vLLM 的 \texttt{CachePolicy} 介面實作，可在 CUDA 與 ROCm 上零程式碼修改部署（ROCm 上的執行期限制見第 \ref{sec:limitations} 節）。
  \item 我們指出既有學習式 KV 管理方法所用的\emph{對稱}損失函數，在動作成本相差 6--190 倍（隨平台而異）的設定下失準（10 倍不對稱下，Bayes 最佳門檻自 $p=0.5$ 移至 $p\approx0.09$），並提出將動作成本納入損失的兩種形式。經比對 8 個既有學習式系統，尚無工作將\emph{異質動作成本}編碼進 KV 管理的損失函數（第 \ref{sec:loss} 節）。
  \item 我們給出完整的實驗計畫，包含以成本項為單位（而非以模組為單位）的消融設計，以及一個離線最優 oracle 作為 go/no-go 判準（第 \ref{sec:eval} 節）。
\end{enumerate}

\paragraph{本稿狀態。}
本文為初稿。第 \ref{sec:motivation} 節的量化分析與第 \ref{sec:related} 節的文獻比對均已完成並經查證；第 \ref{sec:eval} 節為實驗計畫，尚未執行。我們在第 \ref{sec:limitations} 節誠實說明此一狀態的意涵。

% =====================================================================
\section{Background and Motivation}
\label{sec:motivation}

\subsection{KV cache 的記憶體佔用}

對採用 grouped-query attention（GQA）的模型，每 token 的 KV 佔用為
\begin{equation}
  b_{\text{tok}} = 2 \times L \times H_{kv} \times d_{h} \times s
  \label{eq:kvsize}
\end{equation}
其中 $L$ 為層數，$H_{kv}$ 為 KV head 數，$d_h$ 為 head dimension，$s$ 為每元素位元組數，前導的 2 對應 K 與 V。

表 \ref{tab:kvsize} 給出代表模型的實際數值。所有組態均直接取自 HuggingFace \texttt{config.json}。

\paragraph{權重精度是正交的部署參數。}
本文將 LLM 權重視為 frozen，\emph{且不修改其精度}。式 \eqref{eq:actions} 的動作空間描述的是 \textbf{KV cache 的儲存精度}，與模型權重的精度是兩個獨立的維度：

\begin{center}\footnotesize
\setlength{\tabcolsep}{3pt}
\begin{tabular}{@{}ll@{}}
\toprule
維度 & 取值 \\
\midrule
權重精度（\emph{非本文變數}） & BF16 / FP8 / INT8 / AWQ-INT4 \\
KV 狀態（\emph{本文的決策}） & 式 \eqref{eq:actions} 的六個動作 \\
\bottomrule
\end{tabular}
\end{center}

我們因此不假設任何特定權重精度為「主流」。實驗在代表性的部署配置下進行——平台 A 採 AWQ-INT4（記憶體受限的消費級 GPU 之常見配置），平台 B 採 BF16 與 FP8——並將權重精度的影響列為\emph{敏感度分析}（第 \ref{sec:eval} 節），而非主要貢獻。將兩個維度混為一談會使本文同時變成權重量化、KV 卸載、驅逐與重算四個研究問題，模糊真正的貢獻。

\paragraph{關於 BF16 與 FP16。}
本文所檢視的六個模型其 \texttt{config.json} 的 \texttt{torch\_dtype} 均為 \texttt{bfloat16}。二者皆為 2 位元組，故式 \eqref{eq:kvsize} 的算術不變，但本文一律使用 \textbf{BF16} 以與實際發布的權重一致。

\begin{table}[t]
\centering\small
\caption{KV footprint 與原生上下文長度。所有組態均經 HuggingFace \texttt{config.json} 直接驗證（2026-08-29）。\textbf{原生 ctx} 欄是選模的硬條件：低於評測長度者必須施加 RoPE 外推，而外推本身會降低品質，混淆 $\epsilon$ 的量測。$^{M}$ 為 MoE。}
\label{tab:kvsize}
\setlength{\tabcolsep}{2.6pt}
\begin{tabular}{@{}lrrrrr@{}}
\toprule
模型 & $L$ & $H_{kv}$ & $d_h$ & KV/tok & 原生 ctx \\
\midrule
Qwen2.5-7B-Inst-1M & 28 & 4 & 128 & \textbf{56\,Ki}  & \textbf{1{,}010{,}000} \\
Llama-3.1-8B-Inst  & 32 & 8 & 128 & 128\,Ki & 131{,}072 \\
Qwen3-8B           & 36 & 8 & 128 & 144\,Ki & 40{,}960 \\
Qwen3-14B          & 40 & 8 & 128 & 160\,Ki & 40{,}960 \\
Qwen3-30B-A3B$^{M}$ & 48 & 4 & 128 & 96\,Ki  & 40{,}960 \\
Llama-4-Scout$^{M}$ & 48 & 8 & 128 & 192\,Ki & \textbf{10{,}485{,}760} \\
\bottomrule
\end{tabular}
\end{table}

\subsection{觀察一：單請求容量壓力在 MI300X 上幾乎不存在（但在 24\,GB 級加速器上成立）}

表 \ref{tab:fit} 比較同一組模型在 H100（80\,GB）與 MI300X（192\,GB）上的容納能力。此處固定為 \textbf{BF16 權重}以隔離變因（權重精度的影響見上文與敏感度分析），並保留 10\% 記憶體作為 activation 與碎片化餘裕。平台 A 的實際實驗採 AWQ-INT4 權重，其容納能力見第 \ref{sec:twoplatform} 節。

\begin{table}[t]
\centering\small
\caption{單請求 KV footprint 與容納情形（BF16 權重 + BF16 KV）。\checkmark 表示權重加 KV 可置入。$^\S$ 超出原生 \texttt{max\_position\_embeddings}，需 YaRN 外推（見第 \ref{sec:eval} 節）。}
\label{tab:fit}
\begin{tabular}{llrcc}
\toprule
模型 & ctx & KV & H100 & MI300X \\
 & & & 80\,GB & 192\,GB \\
\midrule
\multirow{3}{*}{Llama-3.1-8B}
 & 128K & 15.6\,GB & \checkmark & \checkmark \\
 & 256K & 31.2\,GB & \checkmark & \checkmark \\
 & \textbf{1M}$^\S$ & \textbf{122\,GB} & $\times$ & \textbf{\checkmark} \\
\midrule
\multirow{2}{*}{Qwen2.5-7B-1M}
 & 256K & 13.7\,GB & \checkmark & \checkmark \\
 & 1M   & 53.4\,GB & \checkmark & \checkmark \\
\midrule
Llama-3.1-70B & 128K & 39.1\,GB & $\times$ & $\times$ \\
\bottomrule
\end{tabular}
\end{table}

三個結論值得強調：

\textbf{(a)} 在 MI300X 上，8B--14B 模型於 128K 上下文僅消耗 HBM 的 8--10\%。既有文獻普遍採用的「單請求、上下文遞增」實驗設定，在此硬體上\emph{量測不到記憶體壓力}。相對地，既有的卸載瓶頸量測~\cite{bottlenecks2026} 顯示在 80\,GB 級硬體上 99\% 的延遲花在傳輸，但該結論建立在容量壓力存在的前提上。

\textbf{(b)} Llama-3.1-8B 的 1M token 上下文（122\,GB KV 加 16\,GB 權重）可完整置入單張 MI300X，而 H100 無法。若啟用 MI300X 原生支援的 FP8，Llama-3.1-70B 的 256K 上下文亦可置入。

\textbf{(c)} 在 \emph{MI300X 上}，真實壓力來自三處：\emph{並行度}（Llama-3.1-8B 在 BF16 權重下約 10 個並行 128K 請求即耗盡 HBM，FP8 權重下約 21 個）、\emph{更大的模型}、以及\emph{跨 session 的 KV 保留}（唯一無上界的需求來源）。前二者是靜態的容量問題，第三者才是真正的策略問題。\emph{在記憶體較小的加速器上（如本文的平台 A，24\,GB）則存在第四個來源：單請求上下文本身}（見第 \ref{sec:twoplatform} 節）。

\subsection{觀察二：卸載的相對代價隨硬體世代惡化}

表 \ref{tab:ratio} 給出關鍵比值。

\begin{table}[t]
\centering\small
\caption{記憶體階層比值與計算-傳輸比。規格取自~\cite{cdna3,mi300xperf2025}。主機鏈路為單向理論值；FLOP/byte 為 dense BF16 峰值除以單向鏈路頻寬，代表一位元組穿越鏈路的時間內可 retire 的算術運算量。}
\label{tab:ratio}
\setlength{\tabcolsep}{3.2pt}
\begin{tabular}{lrrrr}
\toprule
 & Mem BW & Host link & \textbf{Mem:link} & FLOP/byte \\
\midrule
RTX 3090$^{\dagger\ddagger}$ & 0.94\,TB/s & 31.5\,GB/s & \textbf{30:1} & 2{,}254 \\
H100 SXM     & 3.35\,TB/s & 64.0\,GB/s & 52:1  & 15{,}459 \\
H200 SXM     & 4.8\,TB/s  & 64.0\,GB/s & 75:1  & 15{,}459 \\
MI300X       & 5.3\,TB/s  & 63.0\,GB/s & \textbf{84:1}  & \textbf{20{,}752} \\
MI355X       & 8.0\,TB/s  & 63.0\,GB/s & \textbf{127:1} & --- \\
GH200        & $\sim$4\,TB/s & 450\,GB/s & $\sim$9:1 & --- \\
\bottomrule
\end{tabular}
\\[2pt]
{\scriptsize $^\dagger$ RTX 3090 使用 GDDR6X 而非 HBM，PCIe 4.0$\times$16。$^\ddagger$ RTX 3090 之值由 $328\ \text{TC} \times 1{,}695\,\text{MHz} \times 128\,\text{FLOP/clk} = 71.2$\,TFLOPS 除以 31.5\,GB/s 得出~\cite{ga102}。GA10x 的 tensor core 於 \emph{FP32 累加}下為 128 FLOP/clk（FP16 累加為 256），BF16 與 FP16-FP32累加同速。此處易誤取 35.6\,TFLOPS——該值對應 64 FLOP/clk，為 TF32 tensor 或非 tensor FP32 的峰值，與 BF16 GEMM 的資料路徑不同。\par}
\end{table}

比值持續惡化的原因是結構性的：HBM 容量與頻寬逐代提升，而 PCIe Gen5 $\times$16 自 MI300X 至 MI355X 完全未變。其後果是，任何卸載策略都必須掩蓋一個比 HBM 存取高近兩個數量級的延遲。

GH200 的 9:1 是例外：藉由 NVLink-C2C 的 450\,GB/s，其卸載經濟性優於任何 x86 主機的加速器約 7--14 倍。\emph{同一套 KV 放置策略不可能同時在 9:1 與 127:1 的平台上最佳}，這是硬體條件化策略的直接論據。

\subsection{觀察三：重算的價值不在單次成本}
\label{sec:recompute-econ}

決定「重算或傳輸」的正確指標並非 HBM 頻寬比，而是 FLOPS 對主機鏈路之比，即一個位元組穿越 PCIe 的時間內，加速器能retire 多少算術運算。MI300X 的 dense BF16 峰值為 1307.4\,TFLOPS，故每傳輸一位元組可換得約 20{,}750 FLOP；H100 為 989.4\,TFLOPS，對應 15{,}460 FLOP。MI300X 在此指標上優於 H100 約 1.34 倍。

然而，在絕對值上傳輸仍然較便宜，且\emph{便宜的程度隨平台劇烈變動}。以簡化的 prefill 模型（線性層約 $2N_{\text{params}}$ FLOP/token，100\% MFU）估算，成本比 $\kappa = t_{\text{recompute}} / t_{\text{transfer}}$ 如表 \ref{tab:kappa}。

\begin{table}[t]
\centering\small
\caption{重算對傳輸的成本比 $\kappa$。$\kappa$ 跨平台變動達 \textbf{32 倍}，這是本文「硬體條件化策略」主張最直接的量化依據。}
\label{tab:kappa}
\begin{tabular}{llrrr}
\toprule
平台 & 模型 & 重算 & 傳輸 & $\boldsymbol{\kappa}$ \\
\midrule
\multirow{2}{*}{MI300X}
 & Llama-3.1-8B & 12.2\,$\mu$s & 2.08\,$\mu$s & \textbf{6}$\times$ \\
 & Llama-3.1-70B & 107\,$\mu$s  & 5.20\,$\mu$s & \textbf{21}$\times$ \\
\midrule
\multirow{2}{*}{RTX 3090}
 & Llama-3.1-8B & 225\,$\mu$s  & 4.16\,$\mu$s & \textbf{54}$\times$ \\
 & Llama-3.1-70B & 1{,}972\,$\mu$s & 10.4\,$\mu$s & \textbf{190}$\times$ \\
\bottomrule
\end{tabular}
\end{table}

且重算成本隨 block 的\emph{絕對位置}增長，而傳輸成本則保持恆定。\emph{表 \ref{tab:kappa} 的數值應視為 $\kappa$ 的下界。}

此處須區分兩個容易混淆的量。整段 prefill 的總成本確為序列長度的二次式（attention 為 $O(L^2)$）；但本文的動作空間作用於\emph{固定大小的 block}，而在位置 $P$ 重算一個大小為 $C$ 的 block，其 attention 需讀取前序的 $P$ 個 token，成本為 $O(C\cdot P)$——\textbf{對 $P$ 為線性}。我們在平台 A 上以 $C=2048$、$P\in\{0,4\text{K},8\text{K},16\text{K},24\text{K}\}$ 實測 Llama-3.1-8B（每點三次取中位數），得

\begin{equation}
C_{\text{recompute}}(P) \;=\; 513.0\;\text{ms} \;+\; 26.9\;\text{ms}\times\frac{P}{1000},
\label{eq:recompute-position}
\end{equation}

線性擬合的最大偏差為 1.6\%，五個量測點皆落於直線上，\emph{未觀察到二次成分}。自 $P=0$ 至 $P=24{,}576$ 成本成長 2.29 倍。此係數依平台與模型而定，但\emph{線性}這個形式來自 $O(C\cdot P)$，與平台無關。將重算視為單一常數會在長上下文下系統性低估其成本。

\paragraph{$\kappa$ 亦隨模型架構變動，且方向可事先計算。}
$\kappa$ 的分子由\emph{啟用}參數量決定，分母由 KV 大小決定，而 MoE 使兩者解耦：其 KV 由 attention 層決定（與同規模 dense 模型相同），但重算僅需啟用的專家。以 Qwen3-30B-A3B 為例（48 層、4 個 KV head、128 個專家中啟用 8 個，啟用參數約 3\,B），其 $\kappa$ 在 RTX 3090 上為 27，\emph{僅為同平台 dense 模型（48--108）的約一半}。

這產生一個可證偽的預測：\textbf{在固定平台上，MoE 模型的最佳 \act{Drop} 觸發門檻應高於 dense 模型}。第 \ref{sec:eval} 節將檢驗之。此軸與硬體軸正交，兩者共同構成本文「成本結構決定策略」主張的雙重檢驗。

圖 \ref{fig:ladder} 揭示了一個關鍵的非對稱性：\emph{重算的價值不在於其單次事件成本，而在於它在被需要之前不佔用任何資源}。因此重算的正確條件是「該 KV block 未來被存取的機率足夠低」——低到節省的容量與頻寬能夠抵償萬一被需要時 6--190 倍的代價。

這個推論的重要後果是：\textbf{未來效用預測不是系統的附加元件，而是決定重算動作正確與否的唯一依據}。這與第 \ref{sec:related} 節的文獻分析指向同一結論。

\begin{figure*}[t]
\centering
\begin{tikzpicture}[
  font=\footnotesize\sffamily,
  >={Stealth[length=2.2mm]},
  tier/.style={draw,rounded corners=3pt,thick,minimum width=4.3cm,minimum height=0.78cm,
               align=center,font=\small\sffamily},
  hdr/.style ={font=\scriptsize\bfseries\sffamily,text=black!70,align=center},
  cell/.style={font=\scriptsize\sffamily,align=center,text=black!80},
  note/.style={font=\scriptsize\sffamily,align=left,text=black!65},
]
% ---- 各層 y 座標 ----
\def\ya{0}    \def\yb{-0.92} \def\yc{-1.84} \def\yd{-2.76} \def\ye{-3.68} \def\yf{-4.60}
\def\xt{3.3}  % 階梯中心

\node[tier,fill=gpucol!20,draw=gpucol]  (t1) at (\xt,\ya) {GPU \;·\; BF16};
\node[tier,fill=gpucol!13,draw=gpucol]  (t2) at (\xt,\yb) {GPU \;·\; FP8};
\node[tier,fill=gpucol!6, draw=gpucol]  (t3) at (\xt,\yc) {GPU \;·\; INT4};
\node[tier,fill=cpucol!14,draw=cpucol]  (t4) at (\xt,\yd) {CPU DRAM \;·\; INT8};
\node[tier,fill=ssdcol!14,draw=ssdcol]  (t5) at (\xt,\ye) {SSD \;·\; compressed};
\node[tier,fill=dropcol!16,draw=dropcol,thick] (t6) at (\xt,\yf) {DROP $\rightarrow$ recompute};

% ---- 左：降級 ----
\draw[->,very thick,black!45] (0.75,\ya+0.30) -- (0.75,\yf-0.30);
\node[rotate=90,anchor=south,font=\scriptsize\sffamily,text=black!55] at (0.42,-2.3) {降級 (demote)};
\foreach \y in {\ya,\yb,\yc,\yd,\ye}
  \draw[->,semithick,black!35] (1.25,\y-0.37) -- (1.25,\y-0.55);

% ---- 右：升級 / restore ----
\draw[->,very thick,orange!65!black] (5.9,\yf-0.30) -- (5.9,\ya+0.30);
\node[rotate=90,anchor=north,font=\scriptsize\sffamily,text=orange!55!black] at (6.24,-2.3) {升級 / restore};
\foreach \y in {\yb,\yc,\yd,\ye,\yf}
  \draw[->,semithick,orange!45!black] (5.42,\y+0.55) -- (5.42,\y+0.37);

% ---- 成本欄 ----
\def\xA{8.0}  \def\xB{11.3} \def\xC{14.9}
\node[hdr] at (\xA,0.95) {平時成本\\[-1pt]\scriptsize\mdseries（未被使用時）};
\node[hdr] at (\xB,0.95) {被需要時的成本\\[-1pt]\scriptsize\mdseries（Llama-3.1-8B, /token）};
\node[hdr] at (\xC,0.95) {說明};

\draw[black!25] (6.9,0.52) -- (17.2,0.52);

\node[cell] at (\xA,\ya) {\textbf{最高}（HBM 最稀缺）};
\node[cell] at (\xA,\yb) {$\tfrac{1}{2}$ of BF16};
\node[cell] at (\xA,\yc) {$\tfrac{1}{4}$ of BF16};
\node[cell] at (\xA,\yd) {低（DRAM）};
\node[cell] at (\xA,\ye) {極低（NVMe）};
\node[cell,text=dropcol!60!black] at (\xA,\yf) {\textbf{零}};

\node[cell] at (\xB,\ya) {0};
\node[cell] at (\xB,\yb) {$\approx$0（反量化）};
\node[cell] at (\xB,\yc) {$\approx$0（反量化）};
\node[cell] at (\xB,\yd) {$\sim$1.0\,$\mu$s};
\node[cell] at (\xB,\ye) {$\sim$9\,$\mu$s};
\node[cell,text=dropcol!60!black] at (\xB,\yf) {\textbf{12.3\,$\mu$s} @100\% MFU};

\node[note,text width=2.5cm] at (\xC,\ya) {最快，但佔用最貴的資源};
\node[note,text width=2.5cm] at (\xC,\yb) {\textbf{生產環境最常見的量化}};
\node[note,text width=2.5cm] at (\xC,\yc) {精度分層};
\node[note,text width=2.5cm] at (\xC,\yd) {位置分層};
\node[note,text width=2.5cm] at (\xC,\ye) {容量近乎無限};
\node[note,text width=2.5cm,text=dropcol!60!black] at (\xC,\yf) {\textbf{與 SSD 同數量級}};

% ---- 關鍵標註 ----
\draw[dropcol,thick,dashed,rounded corners=3pt]
  ($(\xB,\ye)+(-1.5,0.32)$) rectangle ($(\xB,\yf)+(1.5,-0.32)$);
\node[font=\scriptsize\bfseries\sffamily,text=dropcol!55!black,anchor=north]
  at ($(\xB,\yf)+(0,-0.42)$) {平台 A 實測：位置 $>$7K 後 SSD 較重算廉 $\Rightarrow$ 二者應同時決策};

\end{tikzpicture}
\caption{\sys 的動作空間與其成本結構。與既有工作的差異有二：（i）\emph{位置分層}（GPU/CPU/SSD）與\emph{精度分層}（BF16/FP8/INT4/INT8）被統一為單一階梯，而既有工作只做其中之一；（ii）「丟棄後重算」是階梯上的一層，而非檢索失敗後的補救。此六階為典型配置之示意，量化器可替換（見第 \ref{sec:formulation} 節）。關鍵的非對稱性在「平時成本」欄：\textbf{重算在被需要之前不佔用任何資源}，這是它唯一的優勢，也是它唯一的正當理由。右下虛線框標示一個常被忽略的事實：自 SSD 載入與重算處於同一數量級，兩者本就應在同一個決策中被比較。此點已由平台 A 的端到端量測證實，但\emph{量級與歸因均與純頻寬推算不同}——以 Llama-3.1-8B 於 16K 上下文計，自磁碟階取回為 346\,$\mu$s/token、重算為 326\,$\mu$s/token（皆已扣除 GPU 命中的基準），確為同一數量級，然其絕對值較「$\sim$9\,$\mu$s/token @ 7\,GB/s」的理想頻寬推算高出約 30 倍。分解後可見磁碟 I/O 僅佔搬運成本的約 35\%，其餘為分層機制的查表、cascade 排程與 promotion 開銷；將裝置由 SATA SSD（實測 0.49\,GB/s）換為 NVMe（1.65\,GB/s，頻寬 3.4 倍）後端到端成本\emph{未降反升}。詳見第 \ref{sec:eval} 節。}
\label{fig:ladder}
\end{figure*}

\subsection{觀察四：prefill 與 decode 的卸載經濟性相差五個數量級}
\label{sec:prefill-decode}

同一筆搬移在 prefill 攤提於整段前綴，在 decode 則每產生一個 token 就重付一次。這個不對稱決定了動作空間的哪些層在何種工作負載下可用，而既有工作多半只在其中一側成立。

兩個階段的瓶頸不同。prefill 一次處理 $N$ 個 token，算術強度高，受限於算力；decode 一次處理一個 token，但\emph{每一步都須讀取整個 KV cache}，算術強度趨近 1，受限於記憶體頻寬。我們在平台 A 上實測此差距（Llama-3.1-8B 與 Qwen2.5-7B-1M，AWQ 權重，四個上下文長度）：

\begin{center}
\footnotesize
\begin{tabular}{@{}lrrr@{}}
\toprule
上下文 & prefill (ms/tok) & decode (ms/tok) & 倍數 \\
\midrule
16{,}384 & 0.316 & 17.0 & 54$\times$ \\
65{,}536 & 0.591 & 23.5 & 40$\times$ \\
126{,}976 & 0.903 & 29.7 & 33$\times$ \\
258{,}048 & 1.264 & 29.4 & 23$\times$ \\
\bottomrule
\end{tabular}
\end{center}

此差距使卸載的成本結構在兩階段完全不同。以 Llama-3.1-8B 於 126{,}976 上下文、FP8-KV 計，KV 總量為 7.75\,GiB。該平台的 PCIe 為 Gen3 $\times$16（15.75\,GB/s），故將全部 KV 移至主機需 492\,ms。在 prefill，此成本攤提於 126{,}976 個 token，等於每 token 0.0024\,ms；在 decode，\emph{每一步都要重付}，而全駐留於 HBM 的同一次讀取僅需 8.3\,ms。

\begin{center}
\footnotesize
\begin{tabular}{@{}lrr@{}}
\toprule
卸載比例 & decode 每步增加 & 相對全駐留 \\
\midrule
5\% & $+$24.6\,ms & 3.0$\times$ \\
20\% & $+$98.4\,ms & 11.9$\times$ \\
100\% & $+$492.1\,ms & 59.3$\times$ \\
\bottomrule
\end{tabular}
\end{center}

\textbf{在 full attention 下，decode 階段的卸載在任何比例都不可行。} 即使僅卸載 5\%，每步仍付出全駐留 3 倍的代價。唯一可行的路徑是稀疏 attention——僅取回當步 attention 分數高的 block，即 Quest~\cite{quest2024}、InfiniGen~\cite{infinigen2024} 與 ArkVale~\cite{arkvale2024} 的做法。

\paragraph{三種工作負載型態的有效動作空間不同。}

\begin{center}
\footnotesize
\begin{tabular}{@{}p{0.13\columnwidth}p{0.23\columnwidth}p{0.20\columnwidth}p{0.26\columnwidth}@{}}
\toprule
型態 & 決策點 & 預測的對象 & 有效動作空間 \\
\midrule
prefill 為主 & 請求\emph{之間}保留哪些前綴 & 該前綴是否再被請求（時序） & 六階全部 \\
\addlinespace
多輪對話 & 同上，另含思考空窗的機會成本 & 同上 & 六階全部 \\
\addlinespace
生成為主 & 請求\emph{之內}哪些 block 可離開 GPU & 該 block 的 attention 分數（語意） & \textbf{僅 GPU 內的精度階}；CPU/SSD 慢 3--59$\times$ \\
\bottomrule
\end{tabular}
\end{center}

\paragraph{Takeaway.}
第 \ref{sec:eval} 節的實驗涵蓋前兩種型態：其工作負載為「同一前綴送兩次」，prompt 與生成的比例為 1092:1，故 99.8\% 的時間落在 prefill。第三種型態的決策點在請求之內，且預測的對象是 attention 分數而非存取機率——那是不同的預測問題，本文不處理。此界線與觀察三的 \act{SSD} 結果同構：\emph{六階是候選集合，其有效子集由平台與工作負載型態共同決定，而兩者皆可事先算出。}

\subsection{雙平台實驗設計}
\label{sec:twoplatform}

表 \ref{tab:ratio} 的 FLOP/byte 欄位揭示了一個跨越一個數量級的光譜：RTX 3090 為 2{,}254，MI300X 為 20{,}752，相差 9.2 倍。此差距的來源是結構性的：3090 的 GDDR6X 頻寬與 PCIe 4.0 鏈路較為接近，而 MI300X 的 HBM3 遠快於其 PCIe 5.0 鏈路。

這使得本文的核心主張成為可證偽的：\emph{若最佳 KV 放置策略確實隨硬體比值結構性改變，則同一組策略在相距 9.2 倍的此二平台上應展現不同的最佳動作組合}。須說明的是，此二平台位於該光譜的中段而非兩端（表 \ref{tab:ratio} 中 GH200 為 9:1、MI355X 為 127:1），我們的可證偽性主張因而受限於實際可取得的硬體跨度。具體而言，我們預測\act{Drop}$\rightarrow$重算動作在 MI300X 上的最佳觸發門檻，應顯著低於在 RTX 3090 上的門檻。

兩個平台因此承擔不同的實驗職責（表 \ref{tab:platforms}）。值得強調的是，這並非資源限制下的妥協：\emph{單一平台無法檢驗硬體條件化的主張}。

\begin{table}[t]
\centering\small
\caption{雙平台的實驗職責分工。}
\label{tab:platforms}
\begin{tabular}{p{2.3cm}p{2.4cm}p{2.4cm}}
\toprule
 & \textbf{RTX 3090} & \textbf{MI300X} \\
\midrule
記憶體 & 24\,GB GDDR6X & 192\,GB HBM3 \\
壓力來源 & 單請求上下文 & 並行度／跨 session \\
主要壓力軸 & context 8K$\rightarrow$64K & 並行度 1$\rightarrow$32 \\
問題性質 & 容量 & 機會成本 \\
生態 & CUDA（baseline 齊全） & ROCm（部分受限） \\
能耗量測 & 整卡 & \textbf{HBM 分軌} \\
角色 & 機制驗證 + 完整對照 & 規模 + 主要科學主張 \\
\bottomrule
\end{tabular}
\end{table}

\paragraph{RTX 3090 上的容量分析。}
以 8-bit 權重的 Llama-3.1-8B 計算，24\,GB 中約 13.6\,GB 可供 KV 使用，對應約 111K token。單請求在 64K 上下文需 7.8\,GB（總計 15.8\,GB，可置入），而 128K 需 15.6\,GB（總計 23.6\,GB，\emph{超出}）。64K 與 128K 之間存在一道明確的容量懸崖，\emph{既有文獻所預設的場景在此平台上完全成立}。

\paragraph{CUDA 生態的對照組優勢。}
第 \ref{sec:eval} 節指出 ArkVale、Quest、ShadowKV、ClusterKV 與 RetroInfer 因 \texttt{rapidsai/raft} 缺乏 ROCm 移植而無法在 MI300X 上執行。這些系統在 RTX 3090 上均可執行，雙平台設計同時解決了 AMD 平台上對照組不足的問題。

% =====================================================================
\section{Related Work}
\label{sec:related}

我們依動作與訊號兩個維度組織既有工作，並在表 \ref{tab:matrix} 給出特徵矩陣。判準從嚴：「未來預測」欄僅在該工作預測\emph{尚未發生}的存取時才標記，僅使用當前步 attention 分數者不計。

\subsection{分層儲存}
FlexGen~\cite{flexgen2023} 首次以線性規劃決定 KV 在 GPU/CPU/disk 三層間的比例配置，但其策略為離線靜態。CachedAttention~\cite{cachedattention2024} 針對多輪對話保留 session 級 KV 於 HBM/DRAM/SSD，並利用排程器佇列進行預取。Mooncake~\cite{mooncake2025} 將叢集中閒置的 CPU/DRAM/SSD 池化為共享 KVCache 層。KVDrive~\cite{kvdrive2026} 以 prefill 階段的 attention 重要性剖析決定 HBM/DRAM/SSD 的初始配置，並將全量 KV 寫入 SSD 作為後備。Tutti~\cite{tutti2026} 則專注於 SSD I/O 路徑本身，透過 GPU 發起的儲存存取消除 CPU 瓶頸。

這些工作共同的限制是：放置訊號來自當前或過去的觀測，且不將重算視為動作。

\subsection{可恢復驅逐}
ArkVale~\cite{arkvale2024} 觀察到 token 重要性隨解碼步驟變動，早期被驅逐者可能重新變得重要，因而以 page digest（key 向量的包絡體）估計已驅逐頁的重要性並召回。QEvict~\cite{qevict2026} 則將可恢復性建立在\emph{精度}維度：中等信心的視窗被量化保存，重新變重要時反量化升回全精度。

這兩類工作分居兩個正交維度：ArkVale 及其後繼者~\cite{quest2024,shadowkv2025} 做\emph{位置分層}（GPU$\leftrightarrow$CPU，且兩層皆為全精度），QEvict 做\emph{精度分層}（全精度$\leftrightarrow$量化，但僅一個位置）。據我們所知，尚無工作以單一重要性驅動策略統一管理兩者。

\subsection{學習式未來效用預測}
2026 年出現多篇以學習方法預測 KV 未來效用的工作：KVP~\cite{kvp2026} 以 per-head 的強化學習代理排序 token，其 reward 直接導自未來效用；ForesightKV~\cite{foresightkv2026} 建構以真實未來 attention 分數計算的 golden eviction oracle 並以模仿學習逼近；LookaheadKV~\cite{lookaheadkv2026} 則以可學習的 lookahead token 預測模型真實回應所導出的重要性。Marconi~\cite{marconi2025} 的效用函數同時包含預測重用機率、計算節省與記憶體佔用。

這些工作的動作空間\emph{均為單層的保留/丟棄}，且無一支援恢復。\emph{最強的未來預測器沒有安全網，而最完善的安全網沒有未來預測器}。

\subsection{重算與傳輸}
Cake~\cite{cake2025} 讓 GPU 正向計算與 I/O 反向載入同時進行並在中點會合。KVPR~\cite{kvpr2025} 以 profiler 與 scheduler 求解活化值傳輸與 KV 重算的最佳切分。HCache~\cite{hcache2025} 儲存中間 hidden state（約為 KV 的一半大小）並由其重建 KV，實質上在階梯上增加了第三階。

這些工作共同的時序特性是\emph{檢索時的反應式決策}：它們回答「此刻需要這段 KV，如何最快取得」，而非「未來可能需要，現在應置於何處」。ACL 2026 的綜述~\cite{kvsurvey2026} 中，整理「GPU KV 重算 $\leftrightarrow$ KV 傳輸」重疊的表格僅列出一個範例（KVPR）。

\paragraph{以 query 為訊號者無法取代放置決策。}
Quest~\cite{quest2024}、InfiniGen~\cite{infinigen2024}、ArkVale~\cite{arkvale2024} 與 ShadowKV~\cite{shadowkv2025} 皆以\emph{當下的 query 向量}估計各 block 的重要性並取 top-$k$。這是取得單步 attention 稀疏性最直接的訊號，其準確度上限也高於任何以存取歷史為輸入的預測器。然而該類方法在本文的問題上有兩個結構性障礙。

\emph{(i) query 只覆蓋當下這一步。}$q \cdot k$ 回答的是「此刻這個 query 需要哪些 KV」。跨請求的重用——一段前綴在數分鐘後被另一個尚未抵達的請求命中——所對應的 query \emph{尚不存在}，因而無法由此類訊號涵蓋。這正是上一段所述反應式與前瞻式的分界。

\emph{(ii) 計算 $q \cdot k$ 預設 $k$ 仍可取得。}要以 query 判斷是否保留某 block，必須先讀取該 block 的 key；而是否保留該 key 正是待決策的變數。此依賴使 query 訊號無法用於\emph{已降級至 SSD 或已丟棄}的 block，亦即無法用於階梯下緣——而階梯下緣正是成本差異最大之處。以 query 為訊號的方法因此運作於「內容皆可取得，僅需挑選搬運對象」的前提之上，其前一步（內容該置於何處、是否保留）即為本文所處理的問題。

本文不將二者視為競爭關係：query 側訊號可作為\emph{特徵}併入前瞻式預測器（第 \ref{sec:predictor} 節），對仍可取得的 block 提供高品質輸入，而預測器負責覆蓋 query 無法觸及的跨請求與已降級情形。

\subsection{品質作為約束}
KVServe~\cite{kvserve2026} 將壓縮 profile 選擇形式化為
$\arg\min_p T_p(c)$ s.t. $T_p(c) \le T_{\text{SLO}} \wedge q_p(w) \ge q_{\min}$，是我們所知唯一在 Tier-1 場所將模型品質寫為硬性約束的 KV 工作，但其決策變數不含放置。EvicPress~\cite{evicpress2025} 與 AdaptCache~\cite{adaptcache2025} 將品質納入效用函數，但作為可調權重而非約束。LeoAM~\cite{leoam2025} 是我們調查中唯一給出數值上界者（宣稱準確度下降小於 1\%）。OrbitFlow~\cite{orbitflow2026} 以 ILP 於延遲 SLO 與記憶體容量約束下決定逐層放置，但因其為無損設計，ILP 中不存在品質變數。

\subsection{形式化與上界}
Fang 等人~\cite{hetmem2025} 將異質記憶體系統中的 KV 放置問題數學形式化並推導理論上界，明確指出「存在可觀的執行期最佳化空間」，但同樣明確表示\emph{不提出具體排程策略}。本文正是填補其留下的策略空缺。

\subsection{特徵矩陣}

\begin{table*}[t]
\centering\small
\caption{特徵矩陣。$\bullet$ 表示明確支援，$\circ$ 表示部分或間接支援。判準從嚴（見正文）。「未來預測」不含僅使用當前步 attention 估計者。}
\label{tab:matrix}
\begin{tabular}{llccccccc}
\toprule
工作 & 場所 & GPU/CPU & SSD & 未來預測 & 壓縮 & 恢復 & \textbf{重算} & \textbf{品質約束} \\
\midrule
FlexGen~\cite{flexgen2023}        & ICML'23      & $\bullet$ & $\bullet$ &           & $\bullet$ & $\bullet$ &           &           \\
InfiniGen~\cite{infinigen2024}    & OSDI'24      & $\bullet$ &           & $\circ$   &           & $\bullet$ &           &           \\
ArkVale~\cite{arkvale2024}        & NeurIPS'24   & $\bullet$ &           &           &           & $\bullet$ &           &           \\
Cake~\cite{cake2025}              & ICML'25      & $\bullet$ & $\bullet$ &           &           & $\bullet$ & $\bullet$ &           \\
HCache~\cite{hcache2025}          & EuroSys'25   & $\bullet$ & $\bullet$ &           & $\circ$   & $\bullet$ & $\bullet$ &           \\
Marconi~\cite{marconi2025}        & MLSys'25     &           &           & $\bullet$ &           &           & $\circ$   &           \\
ShadowKV~\cite{shadowkv2025}      & ICML'25      & $\bullet$ &           &           & $\bullet$ & $\bullet$ & $\circ$   &           \\
OrbitFlow~\cite{orbitflow2026}    & VLDB'26      & $\bullet$ &           &           &           & $\bullet$ &           &           \\
KVServe~\cite{kvserve2026}        & SIGCOMM'26   & $\bullet$ &           &           & $\bullet$ &           &           & $\bullet$ \\
KVP~\cite{kvp2026}                & ICML'26      &           &           & $\bullet$ &           &           &           &           \\
Het-Mem~\cite{hetmem2025}         & IEEE CAL'25  & $\bullet$ &           &           &           & $\bullet$ &           &           \\
MTDS~\cite{mtds2026}              & C\&IS'26     & $\bullet$ & $\bullet$ & $\circ$   &           & $\bullet$ & $\bullet$ &           \\
KVDrive~\cite{kvdrive2026}        & preprint     & $\bullet$ & $\bullet$ &           &           & $\bullet$ &           &           \\
EvicPress~\cite{evicpress2025}    & preprint     & $\bullet$ & $\bullet$ & $\circ$   & $\bullet$ & $\bullet$ &           & $\circ$   \\
QEvict~\cite{qevict2026}          & preprint     &           &           &           & $\bullet$ & $\bullet$ &           &           \\
\midrule
\textbf{\sys（本文）}              & ---          & $\bullet$ & $\bullet$ & $\bullet$ & $\bullet$ & $\bullet$ & $\bullet$ & $\bullet$ \\
\bottomrule
\end{tabular}
\end{table*}

\subsection{研究缺口}

由表 \ref{tab:matrix} 可得三項觀察：

\textbf{(1) 「重算」與「品質約束」兩欄從未在同一列同時標記。}

\textbf{(2) 四項成本從未同時進入單一線上決策。}傳輸成本見於 CacheGen 一系，重算成本見於 Cake/KVPR/HCache，記憶體成本見於 EvicPress，未來效用見於 Marconi/KVP。最高覆蓋為 Marconi 的三項（未來效用、計算節省、記憶體佔用），其因僅在 GPU 上運作而\emph{結構性地}不可能具備傳輸成本。

\textbf{(3) 兩條研究血脈互不交會。}以 FlexGen 為基礎的一系（IMPRESS、LeoAM、InfiniGen）具備重要性感知與有損處理，但無跨請求重用；以 vLLM 為基礎的一系（CacheBlend~\cite{cacheblend2025}、LMCache~\cite{lmcache2025}、Cake、MTDS）具備跨請求重用，但將 KV 視為精確且不可分割。\emph{尚無系統同時具備重要性感知的有損分層與跨請求重用}。

綜述~\cite{kvsurvey2026} 自身的 Observation O7 指出壓縮與放置的協同設計是「a missed opportunity」，Challenge C5 呼籲在共享預算下聯合最佳化 eviction、offload 與 prefetch，而該清單\emph{未包含重算}。

% =====================================================================
\section{Problem Formulation}
\label{sec:formulation}

\subsection{狀態與動作}

設 KV cache 被劃分為 block 集合 $\mathcal{B}$。於決策時刻 $t$，系統狀態為
\begin{equation}
  s_t = \Big( \{ \ell_i, \pi_i, a_i, n_i \}_{i \in \mathcal{B}},\; C^{\text{gpu}}, C^{\text{cpu}}, C^{\text{ssd}},\; Q_{\text{pcie}},\; \phi_t \Big)
\end{equation}
其中 $\ell_i$ 為 block $i$ 的所在層級，$\pi_i$ 為其精度，$a_i$ 為上次存取時刻，$n_i$ 為 token 數，$C^{\ast}$ 為各層剩餘容量，$Q_{\text{pcie}}$ 為互連佇列深度，$\phi_t$ 為當前查詢特徵。

動作空間為
\begin{equation}
  \mathcal{A} = \{\, \act{Gpu}_{16},\; \act{Gpu}_{8},\; \act{Gpu}_{4},\; \act{Cpu}_{8},\; \act{Ssd},\; \act{Drop} \,\}
  \label{eq:actions}
\end{equation}
對應圖 \ref{fig:ladder}。\act{Drop} 動作意味著該 block 不被儲存於任何位置，若後續被需要則由 token 重算。

\paragraph{階梯是示意，量化器可替換。}
式 \eqref{eq:actions} 的六個狀態代表\emph{典型}部署選項，而非本文主張的特定量化方案。實務上 KV 量化存在一個重要的非對稱性：\textbf{key 對量化誤差遠比 value 敏感}，故生產系統普遍採 K/V 分離的精度配置——LMCache 的預設 \texttt{turboquant\_k8v4} 為 key 8 位元、value 4 位元；KIVI~\cite{kivi2024} 對 key 採 per-channel、對 value 採 per-token 量化，因二者離群值結構不同。

本文的決策框架與此正交：任何量化器只要能給出「容量佔用」與「存取延遲」兩個數字，即可作為階梯的一階。將 \act{Gpu}$_{8}$ 替換為 K8V4、或將 \act{Gpu}$_{4}$ 替換為 KIVI 的 2 位元方案，均不改變式 \eqref{eq:opt} 的形式，只改變成本常數。

\paragraph{\act{Drop} 的可行性受前序 block 約束。}
其餘五個動作彼此獨立——一個 block 放在哪一層，不影響其他 block 能否被取回。\textbf{\act{Drop} 是唯一的例外。}

重算 block $i$ 並不需要自 token 0 重跑整個序列：以 chunked prefill 的方式，只需將 block $i$ 的 token 前向通過模型，其 attention 讀取\emph{前序 block 既有的 KV}。這使重算成本為 $O(|b_i|)$ 而非 $O(i)$。但這也意味著：

\begin{equation}
  a_i = \act{Drop}\ \text{可行} \iff \forall j < i:\ \text{block } j \text{ 的 KV 可取得}
  \label{eq:dropdep}
\end{equation}

若前序 block 亦被丟棄，必須先重算它們，成本沿依賴鏈累積：
\begin{equation}
\begin{aligned}
  \mathrm{recomp}(i) &= \sum_{j \in \mathcal{D}(i)} \Big( 2 N_{\text{active}} |b_j| + \mathrm{attn}(j) \Big) \\
  \mathcal{D}(i) &= \{i\} \cup \{\, j < i : a_j = \act{Drop} \,\}
\end{aligned}
\end{equation}

\textbf{此約束的意涵是 \act{Drop} 決策彼此耦合，不可獨立最佳化。} 一個只看單一 block 效用的策略會系統性地低估連續丟棄的代價。我們以「同一序列內連續 \act{Drop} 的 block 數上限」作為可行性約束處理之，並於第 \ref{sec:eval} 節量測該上限對 Pareto 前緣的影響。這也是 HCache~\cite{hcache2025} 儲存 hidden state 的動機：由 hidden state 可單步重建 KV，從而\emph{切斷依賴鏈}。

\subsection{成本模型}

對 block $i$ 採取動作 $a$ 的即時成本為
\begin{align}
  c(i,a) ={}& \underbrace{\lambda_{m} \cdot \mathrm{mem}(i,a)}_{\text{容量佔用}}
             + \underbrace{\lambda_{x} \cdot \mathrm{xfer}(i,a)}_{\text{搬移成本}} \nonumber\\
            &+ \underbrace{p_i \cdot \big[\, \mathrm{fetch}(i,a) + \mathrm{recomp}(i,a) \,\big]}_{\text{期望存取成本}}
  \label{eq:cost}
\end{align}
其中 $p_i = \Pr[\text{block } i \text{ 於未來視窗內被存取}]$ 為\emph{未來效用預測}，$\mathrm{fetch}$ 為自 $\ell_i$ 載回 GPU 的延遲，$\mathrm{recomp}$ 僅在 $a = \act{Drop}$ 時非零，且依第 \ref{sec:recompute-econ} 節之模型隨 block 的絕對位置增長。

\subsection{受約束的最佳化問題}

給定品質預算 $\epsilon$，
\begin{equation}
\begin{aligned}
  \min_{\{a_i\}} \quad & \sum_{i \in \mathcal{B}} c(i, a_i) \\
  \text{s.t.} \quad & \Delta\mathcal{Q}\big(\{a_i\}\big) \le \epsilon \\
                    & \textstyle\sum_i \mathrm{mem}^{\text{gpu}}(i,a_i) \le C^{\text{gpu}} \\
                    & \text{(對 CPU/SSD 同理)} \\
                    & \mathrm{TTFT} \le T_{\text{SLO}},\;\; \mathrm{TPOT} \le P_{\text{SLO}} \\
                    & \text{式 \eqref{eq:dropdep} 的可行性約束}
\end{aligned}
\label{eq:opt}
\end{equation}

此形式化與既有工作的差異有三：（i）動作空間 \eqref{eq:actions} 同時跨越位置與精度兩個維度，並將 \act{Drop} 視為一個具明確標價的層級而非失敗路徑；（ii）品質為\emph{硬性約束}而非目標函數中的加權項，此點與 KVServe~\cite{kvserve2026} 一致但作用於放置而非壓縮；（iii）成本函數 \eqref{eq:cost} 中四項成本同時出現。

\paragraph{$\epsilon$ 的可操作性。}
$\Delta\mathcal{Q}$ 於執行期不可直接觀測，這是所有品質約束工作面臨的共同困難。KVServe 採用離線 profiling 加上 $\epsilon$-greedy bandit 修正。我們規劃在第 \ref{sec:eval} 節評估兩種代理訊號：（a）離線 profile 建立的 (壓縮率, 任務類型) $\rightarrow$ 品質查找表；（b）模型輸出分佈的信心值作為線上回饋。此為本設計最主要的風險點，我們於第 \ref{sec:limitations} 節說明。

% =====================================================================
\section{Design}
\label{sec:design}

\subsection{系統架構}

\begin{figure*}[t]
\centering
\begin{tikzpicture}[
  font=\footnotesize\sffamily,
  >={Stealth[length=2.2mm]},
  eng/.style  ={draw,rounded corners=3pt,fill=black!4,draw=black!45,thick,align=center},
  sig/.style  ={draw,rounded corners=2pt,fill=white,draw=black!35,font=\scriptsize\sffamily,
                align=center,minimum height=0.58cm,minimum width=3.15cm,inner sep=2.5pt},
  stage/.style={draw,rounded corners=3pt,fill=white,draw=orange!65!black,thick,align=center,
                minimum height=1.15cm,minimum width=2.35cm},
  tier/.style ={draw,rounded corners=3pt,thick,align=center,minimum height=0.74cm,
                minimum width=2.75cm,font=\scriptsize\sffamily},
  ar/.style   ={->,thick,black!62},
  tar/.style  ={->,semithick,black!40},
]

% ================= Row 1 : 既有引擎 =================
\node[eng,minimum width=3.8cm,minimum height=0.85cm] (llm)  at (3.0,0)
      {Frozen Llama / Qwen\\[-2pt]\scriptsize Attention};
\node[eng,minimum width=5.4cm,minimum height=0.85cm] (vllm) at (9.1,0)
      {vLLM \;\textbf{·}\; \texttt{OffloadingConnector}\\[-2pt]\scriptsize PagedAttention · prefix cache};
\draw[ar] (llm) -- (vllm);

% ================= Row 2 : 訊號 + Tiara =================
\node[sig] (s1) at (2.78,-1.28) {block 狀態 $\ell_i,\pi_i,a_i,n_i$};
\node[sig] (s2) at (2.78,-1.90) {query 特徵 $\phi_t$};
\node[sig] (s3) at (2.78,-2.52) {剩餘容量 $C^{\text{gpu/cpu/ssd}}$};
\node[sig] (s4) at (2.78,-3.14) {互連佇列深度 $Q_{\text{pcie}}$};
\node[font=\scriptsize\itshape,text=black!55,anchor=south] at (2.78,-1.02) {執行期訊號};

\node[stage] (pred)  at (6.30,-2.15) {未來效用預測器\\[2pt]$p_i=\Pr[\text{未來被存取}]$};
\node[stage] (cost)  at (9.05,-2.15) {成本模型\\[2pt]式~\eqref{eq:cost}};
\node[stage] (qgate) at (11.80,-2.15) {品質閘\\[2pt]$\Delta\mathcal{Q}\le\epsilon$};
\draw[ar] (pred) -- (cost);
\draw[ar] (cost) -- (qgate);

\begin{scope}[on background layer]
  \node[draw=orange!70!black,thick,rounded corners=4pt,fill=hlbox,
        fit=(pred)(cost)(qgate),inner sep=9pt] (tiara) {};
\end{scope}
\node[anchor=south west,font=\scriptsize\bfseries\sffamily,text=orange!55!black]
  at ($(tiara.north west)+(0.05,0.04)$) {\sys \;— 本文貢獻};

\foreach \n in {s1,s2,s3,s4} \draw[tar] (\n.east) -- (tiara.west);

% CachePolicy hook：從 vLLM 垂直下到 Tiara
\draw[ar,dashed,black!55] (vllm.south) -- (tiara.north)
  node[midway,left=3pt,font=\scriptsize\itshape,text=black!55,align=right]
  {掛載於\\[-1pt]\texttt{CachePolicy}};

% 動作輸出
\node[anchor=west,align=left,font=\scriptsize\sffamily] (actlbl) at (14.35,-2.15)
  {動作 $a_i\in\mathcal{A}$\\[-1pt]\scriptsize 式~\eqref{eq:actions}};
\draw[ar] (tiara.east) -- (actlbl.west);

% ================= Row 3 : Memory Manager（匯流排） =================
\node[eng,minimum width=15.0cm,minimum height=0.62cm] (mm) at (8.72,-4.15)
      {Memory Manager \;\scriptsize（執行搬移 / 壓縮 / 重算）};
\draw[ar] (tiara.south) -- (tiara.south|-mm.north);

% ================= Row 4 : 六個層級 =================
\node[tier,fill=gpucol!18,draw=gpucol]  (t1) at (2.58,-5.45) {GPU · BF16};
\node[tier,fill=gpucol!8, draw=gpucol]  (t2) at (5.65,-5.45) {GPU · INT4};
\node[tier,fill=cpucol!14,draw=cpucol]  (t3) at (8.72,-5.45) {CPU · INT8};
\node[tier,fill=ssdcol!14,draw=ssdcol]  (t4) at (11.79,-5.45) {SSD · 壓縮};
\node[tier,fill=dropcol!16,draw=dropcol](t5) at (14.86,-5.45) {DROP\\[-3pt]$\rightarrow$ recompute};

\foreach \t in {t1,t2,t3,t4,t5} \draw[tar] (\t.north|-mm.south) -- (\t.north);

% ================= restore：最左側獨立通道 =================
\coordinate (rl) at (0.42,-5.45);
\draw[ar,dashed,orange!60!black]
  (t1.west) -- (rl) -- (0.42,0) -- (llm.west);
\node[rotate=90,anchor=south,font=\scriptsize\itshape,text=orange!55!black] at (0.10,-2.7)
  {restore / recompute};

\end{tikzpicture}
\caption{\sys 架構與資料流。灰底為既有基礎設施，橘框為本文貢獻。控制器接收四類執行期訊號，經「未來效用預測 $\rightarrow$ 成本模型 $\rightarrow$ 品質閘」三階段，為每個 KV block 產生一個動作 $a_i\in\mathcal{A}$，交由 Memory Manager 執行。關鍵設計決定是\emph{不自行實作推論引擎}：\sys 掛載於 vLLM 既有的 \texttt{CachePolicy} 介面，貢獻邊界因而清晰（\textbf{機制屬於 vLLM，策略屬於本文}），且可在 CUDA 與 ROCm 上零移植部署，並直接與內建的 \texttt{lru}/\texttt{arc} 對照。左側橘色虛線為恢復路徑：被降級或丟棄的 block 在重新變得重要時可升回上層或重算，單次誤判不會造成資訊的永久損失。}
\label{fig:arch}
\end{figure*}

圖 \ref{fig:arch} 給出整體架構。設計上的關鍵決定是：\emph{不自行實作推論引擎}。vLLM~\cite{vllm2023} 的 \texttt{OffloadingConnector} 已提供 GPU$\rightarrow$CPU$\rightarrow$（FS/物件儲存/P2P）的搬移機制，且官方文件明確載明支援 CUDA、ROCm 與 XPU；其上游 CI 於實體 gfx942 硬體上執行 KV offload 測試。該連接器內建的置換策略僅有 \texttt{lru} 與 \texttt{arc}，並透過 \texttt{CachePolicy} 介面開放自訂實作。

此決定有三個後果：（i）\sys 可在 CUDA 與 ROCm 兩個平台上零移植部署；（ii）對照組為 production 級實作而非我們自行撰寫的弱化版本；（iii）貢獻邊界清晰：\emph{機制屬於 vLLM，策略屬於本文}。

\subsection{未來效用預測器}
\label{sec:predictor}

預測器輸出 $p_i$，即 block $i$ 於未來視窗內\emph{被後續請求命中}的機率。此定義的範圍須明確：$p_i$ 是跨請求的時序量，而非請求之內的 attention 分數。在 full attention 的 decode 階段，每個 block 於每一步皆被讀取，$p_i$ 恆為 1，預測無從施力（第 \ref{sec:prefill-decode} 節）。請求之內的取捨由稀疏 attention 處理，其預測對象為 $q \cdot k$，屬不同問題。我們對 8 個既有的學習式快取／放置系統做了設計層級的比對（表 \ref{tab:predictors}），並據此做出四項設計決定。

\begin{table}[t]
\centering\small
\setlength{\tabcolsep}{3pt}
\caption{既有學習式快取／放置預測器的設計比對。最後一欄為本文關注的重點：損失函數是否對不同錯誤方向給予不同代價。}
\label{tab:predictors}
\begin{tabular}{llcl}
\toprule
系統 & 模型 & 動作數 & 成本不對稱 \\
\midrule
KVP~\cite{kvp2026}          & 112$\times$MLP & 2 & 否 \\
ForesightKV~\cite{foresightkv2026} & MLP(16) & 2 & 部分（單邊） \\
LookaheadKV~\cite{lookaheadkv2026} & LoRA $r{=}8$ & 2 & 否 \\
Marconi~\cite{marconi2025}  & 啟發式 & 2 & 否（但有 value/byte） \\
Mockingjay~\cite{mockingjay2022} & 表格+TD & 2 & 否 \\
LRB~\cite{lrb2020}          & GBDT & 2 & 否（L2） \\
\textbf{PARROT}~\cite{parrot2020} & LSTM+attn & 2 & \textbf{是（NDCG gain）} \\
\textbf{Sibyl}~\cite{sibyl2022} & C51 DQN & $n$ & \textbf{是（實測延遲）} \\
\midrule
\textbf{\sys（本文）} & GBDT & \textbf{6} & \textbf{是（動作成本加權）} \\
\bottomrule
\end{tabular}
\end{table}

\paragraph{規格。}
預測器是一個獨立於 LLM 之外的\emph{小型模型}，不參與 attention 計算，亦不修改模型權重。其介面如下：

\vspace{2pt}
\noindent\fbox{\parbox{0.955\columnwidth}{\footnotesize\sffamily
\textbf{輸入}（每個 KV block，全部於服務時免費取得）\\[2pt]
\hspace*{6pt}$\bullet$ \texttt{deltas[0..31]}：前 32 次存取的時間間隔（未填滿為 $\infty$）\\
\hspace*{6pt}$\bullet$ \texttt{edc[0..9]}：指數衰減計數器，僅於存取時更新\\
\hspace*{6pt}$\bullet$ \texttt{block\_size}, \texttt{layer\_id}, \texttt{head\_id}, \texttt{token\_type}\\
\hspace*{6pt}$\bullet$ block 內 key/value 向量的池化統計\\
\hspace*{6pt}$\bullet$ \texttt{attn\_mass[0..7]}：最近 8 個解碼步中，該 block 取得的\\
\hspace*{20pt}attention 質量佔比（forward 已算出，記錄成本近零）\\[4pt]
\textbf{輸出}\\[2pt]
\hspace*{6pt}$\hat{y}_i = \log \tau_i$ \quad（迴歸；$\tau_i$ 為距下次存取的時間）\\
\hspace*{14pt}$\downarrow$ isotonic regression（驗證集上校準）\\
\hspace*{6pt}$\hat{p}_i = \Pr[\tau_i \le W]$ \quad（$W$ 為決策視窗長度）\\[4pt]
\textbf{模型}：GBDT，單一全域模型，\texttt{layer\_id}/\texttt{head\_id} 作為特徵
}}
\vspace{3pt}

\emph{門檻永遠施加於 $\hat{p}_i$，不施加於 $\hat{y}_i$}——式 \eqref{eq:cost} 與式 \eqref{eq:threshold} 使用的皆為機率。此區分是必要的：對 $\hat{y}_i$ 直接閾值化在數學上不對應式 \eqref{eq:threshold} 的最佳邊界（見第 \ref{sec:loss} 節）。

\paragraph{決定一：預測連續量，而非二元分類。}
所有預測連續量並延後閾值化的系統，都勝過其二元分類的前身，且 Mockingjay~\cite{mockingjay2022} 與 LRB~\cite{lrb2020} 均明確報告此點。LRB 嘗試過對 Belady 邊界的二元分類並依 good-decision-ratio 拒絕之。我們因此回歸 $\log(\text{距下次存取的時間})$。取對數的理由與 LRB 相同：真正重要的是落在邊界的哪一側，100K 與 2M 的差異遠比 2M 與 3M 重要。

\paragraph{決定二：特徵取自 LRB 的三個家族，全部在服務時免費。}
(i) \emph{deltas}，即距該 block 前 $k$ 次存取的時間間隔（$k \le 32$，未填滿者為 $\infty$），此設計涵蓋 LRU、LRU-$K$ 與 S4LRU；(ii) \emph{EDC}，即 10 個指數衰減計數器 $C_i \leftarrow 1 + C_i \cdot 2^{-\Delta_1 / 2^{9+i}}$，僅在存取時更新；(iii) \emph{靜態特徵}，即 block 大小、layer id、head id 與 \emph{token 類型}。最後一項的依據是 SAECache~\cite{saecache2026} 量到不同語意類型（system prompt / 使用者查詢 / 工具輸出 / 推理鏈）的重用率差異達 \textbf{756 倍}，而該特徵取得成本為零。

\paragraph{決定二之補充：必須納入模型內部表徵。}
上述特徵全部來自\emph{存取歷史}，不含模型自身的表徵。這與三個學習式先例形成對比：KVP~\cite{kvp2026} 明確「僅使用 key 與 value 向量」，LookaheadKV~\cite{lookaheadkv2026} 透過層內 adapter 讀取，ForesightKV~\cite{foresightkv2026} 蒸餾真實的未來 attention。\emph{「在不看 KV 的情況下預測 KV 的重用」作為服務層的低成本選擇是合理的，但作為學習貢獻則站不住腳。}我們因此將 block 內 key/value 向量的池化統計併入特徵集，並在第 \ref{sec:eval} 節以消融比較「僅存取歷史」與「存取歷史 + 模型內部表徵」兩種特徵集。

\paragraph{決定二之補充二：納入 query 側訊號，但作為特徵而非機制。}
第 \ref{sec:related} 節論證以 query 為訊號者無法取代放置決策，原因是 $q \cdot k$ 對已降級的 block 不可得。但對\emph{仍位於 GPU} 的 block，該訊號的品質高於任何存取歷史特徵，捨棄它並無理由。我們因此記錄 \texttt{attn\_mass}，即最近 8 個解碼步中該 block 取得的 attention 質量佔總質量之比。此量在 forward pass 中已被計算，記錄它僅需一次 per-block 的規約，不引入額外的 $q \cdot k$。

此設計使特徵集在階梯上呈\emph{降級退化}的形態：位於 GPU 的 block 同時具備存取歷史與 query 側訊號，位於 CPU/SSD 者僅有存取歷史，已丟棄者僅有靜態特徵。GBDT 原生處理缺失值，此結構無須額外處理。我們在第 \ref{sec:eval} 節以消融量測 \texttt{attn\_mass} 的邊際貢獻。

\paragraph{決定三：GBDT 而非神經網路。}
LRB 明確比較過 GBM、線性迴歸、邏輯迴歸、SVM 與 2 層 16 隱藏單元的 NN，GBM 在不同 trace 與快取大小下穩定勝出，且訓練僅需 300\,ms、對 64 個候選推論僅需 \textbf{30\,$\mu$s}。GBDT 亦能原生處理 $\infty$ 值的 delta 特徵。我們採用單一全域模型並以 layer/head id 作為特徵；KVP~\cite{kvp2026} 的 112 個獨立 agent 帶來可觀的工程成本，而其論文並未提供「per-head 模型優於帶 head-id 特徵的共享模型」的消融證據。

此決定另有兩個本設定特有的理由。\emph{(i) 預測器的成本自其收益中扣除。}預測器運行於服務熱路徑，其延遲直接抵銷所避免的重算或傳輸。以表 \ref{tab:kappa} 的區間為參照，單次重算為 12--1{,}972\,$\mu$s，而對 64 個候選的 GBDT 推論為 30\,$\mu$s；若改以需要一次 forward pass 的模型評分，僅 kernel launch 與排隊即進入相同數量級，收益隨即被抵銷。\emph{(ii) 神經預測器與主模型競爭同一加速器。}GBDT 於 CPU 上執行，完全不佔用 GPU；任何在 GPU 上評分的預測器都會插入主模型的 decode 批次，其代價不僅是預測器本身的延遲，還包括對批次結構的擾動。

此外，主模型規模上升時此權衡\emph{朝有利於 GBDT 的方向移動}：$\kappa$ 隨模型增大而上升（表 \ref{tab:kappa}：8B 至 70B 由 6 增至 21，於平台 A 由 54 增至 190），而 GBDT 的推論成本不隨模型規模改變，因其特徵取自存取歷史而非模型參數。

我們不要求讀者接受此論證，而以實驗回答：表 \ref{tab:ablation} 的 (D) 區塊於相同特徵集與相同損失下比較四種預測器，並同時量測\emph{品質收益與熱路徑延遲}。

前三者（GBDT、小型 MLP、其容量放大一階者）皆接受攤平後的特徵向量，故僅檢驗\emph{容量}。第四者為直接消費 \texttt{deltas} 時間序列的小型 GRU，檢驗的是\emph{結構}——即時序資訊應由手工特徵編碼（本文與 LRB 的作法），抑或應由模型自行學習。此對照不可省略：LRB 所比較的 2 層 16 隱藏單元 NN 同樣接受攤平輸入，其「GBM 勝出」的結論因而\emph{不涵蓋}序列模型，本文若僅援引該結論即屬過度推論。

\texttt{deltas} 與 EDC 事實上已是時序的手工編碼——前者為存取間隔序列的截斷，後者為 512 至 262{,}144 十種半衰期上的指數加權和，其功能與循環網路的隱藏狀態同構，差別在於前者的更新規則寫死而後者由訓練得到。(D) 區塊所檢驗的正是此差別是否重要。若 GRU 在計入自身延遲後仍佔優，第 \ref{sec:predictor} 節的整條特徵工程路線即被推翻，我們將據以修改。

\paragraph{決定四：標籤取得採用 LRB 的雙機制。}
標籤 $\tau_i$ 無法在記錄特徵的當下取得，因其定義涉及未來。我們於取樣時記錄特徵並將該樣本置入待標記佇列，其後由兩個機制之一完成標記：\textbf{(a)} 該 block 下次被存取時，以兩次存取的間隔為標籤；或 \textbf{(b)} 當該樣本離開一個長於 relaxed-Belady 邊界的滑動視窗 $W$ 而仍未被存取時，立即以「$>W$」標記。

機制 (b) 是取得\emph{負樣本}的唯一途徑。僅有機制 (a) 時，訓練集只會包含「後來確實被重用」的 block——從未被重用者永遠等不到第二次存取，因而永遠不進入訓練集。模型將完全未見過應被丟棄的樣本長何模樣，而那正是 \act{Drop} 動作的目標群體。取樣以 block 為單位而非以請求為單位，以避免偏向熱門 block。

\paragraph{決定五：以時間切分而非隨機切分評估。}
訓練與測試集依 trace 的\emph{時間順序}切分（前 70\% 訓練、後 30\% 測試），不作隨機切分。隨機切分會使訓練集含有測試期之後的存取事件，而 EDC 與 delta 特徵本身即為時間的函數，此洩漏將系統性地高估預測品質。此點在快取預測文獻中屬常識，但其違反不易由結果察覺，我們因此明確記載。

\paragraph{決定六：以線上重訓因應工作負載飄移。}
不同工作負載的重用結構不同：多輪對話的前綴重複命中，RAG 的文件塊多為一次性讀取，agent 流量則呈現極熱的系統提示與極冷的工具輸出。SAECache~\cite{saecache2026} 所量到的 756 倍重用率差異即為此分布差異的直接證據。我們以三個層次因應：(i) \texttt{token\_type} 作為特徵，使模型得以在單一模型內區分語意類型；(ii) \emph{週期性重訓}，利用 GBDT 300\,ms 的訓練成本，以最近一段時窗的樣本重新擬合；(iii) 第 \ref{sec:eval} 節報告\emph{跨工作負載泛化}，即於工作負載 A 訓練並於 B 測試的退化幅度。

第 (ii) 項是 GBDT 選擇的直接後果而非附帶效果。KVP~\cite{kvp2026} 的 112 個強化學習 agent 需 8 GPU DDP 訓練，其重訓週期以日計，結構上無法追隨以分鐘計的流量變化。\emph{對分布飄移的適應成本，是預測器族選擇的一部分，而非可事後補救的工程細節。}

\subsection{成本敏感的損失函數}
\label{sec:loss}

這是本文與既有工作差異最大的一點，也是我們認為最具貢獻的設計。

\paragraph{問題：對稱損失在本設定下可證明地失準。}
既有 KV cache 的學習式方法（KVP、LookaheadKV、SP-KV）與經典快取方法（LRB、Mockingjay）的損失函數對兩種錯誤方向給予相同代價。在單一動作（驅逐）的設定下這是合理的。但本文的動作空間有六個成本迥異的選項，兩種錯誤的下游代價相差 6--190 倍，且\emph{該比值本身隨平台變動達 32 倍}（表 \ref{tab:kappa}）：

\begin{itemize}\itemsep0.1em
  \item \textbf{False negative}（預測不會被使用 $\rightarrow$ \act{Drop} $\rightarrow$ 後續被需要）：被迫重算，12--108\,$\mu$s/token
  \item \textbf{False positive}（預測會被使用 $\rightarrow$ 保留 $\rightarrow$ 實際未使用）：僅浪費容量
\end{itemize}

令 $\kappa = c_{\text{FN}} / c_{\text{FP}}$ 為成本比。保留的期望成本為 $(1-p)\,c_{\text{FP}}$，丟棄的期望成本為 $p\,c_{\text{FN}}$，兩者相等處即為最佳決策門檻：
\begin{equation}
  p^{*} = \frac{c_{\text{FP}}}{c_{\text{FP}} + c_{\text{FN}}} = \frac{1}{1+\kappa}
  \label{eq:threshold}
\end{equation}
這是成本敏感學習的標準結果~\cite{elkan2001}，我們不主張其為本文貢獻。其意涵是：在 $\kappa = 10$ 時 $p^{*} = 1/11 \approx 0.09$。\emph{一個以對稱損失訓練、並在 $0.5$ 處閾值化的預測器，其決策門檻被錯置了約一個數量級}。

本文的主張是更具體的一點：\textbf{既有所有學習式 KV 管理方法的動作空間都是二元且成本同質的，因此 $\kappa = 1$、$p^{*} = 0.5$，對稱損失恰好正確；而本文的動作空間成本異質且 $\kappa$ 由硬體決定，對稱損失因而失準。}

\paragraph{做法。}
我們採用兩種形式，並在第 \ref{sec:eval} 節比較：

\textbf{(L1) 動作成本加權的 L2}。在 $\log(\text{距下次存取時間})$ 的迴歸上，令每個樣本的權重等於「策略對該 block 所採取的動作，在判斷錯誤時的代價」：
\begin{equation}
  \mathcal{L}_{\text{L1}} = \sum_i w(a_i) \cdot \big( \hat{y}_i - \log \tau_i \big)^2,
  \quad w(a_i) = \frac{\mathrm{cost}_{\text{err}}(a_i)}{\mathrm{cost}_{\text{err}}(\act{Cpu}_8)}
\end{equation}
此形式是 LRB 的損失加上其不需要的那一項，在 LightGBM 中僅需 \texttt{sample\_weight}，無額外機制。

\textbf{但必須注意一個校準上的限制。} 對連續目標的迴歸加權，改變的是條件期望 $\mathbb{E}[\log \tau \mid x]$，\emph{而非決策門檻}。(L1) 本身因而\textbf{不保證}產生式 \eqref{eq:threshold} 的最佳邊界。我們將 (L1) 定位為啟發式，並在其後串接一個明確的門檻校準步驟：於驗證集上將預測值映射為 $\hat{p}$（例如以 isotonic regression），再依式 \eqref{eq:threshold} 設定各層級的門檻。

\paragraph{兩階段做法是正確的基線，本文須證明加權訓練帶來額外增益。}
上述步驟指向一個明顯的替代方案：\emph{以對稱損失訓練，僅在決策時將門檻自 $0.5$ 移至 $p^{*}$}。此即 Elkan~\cite{elkan2001} 的結果，且具備一個實際優點——$\kappa$ 隨平台改變時無需重新訓練。若本文的增益全部來自門檻移動，則加權訓練並無貢獻。我們因此在表 \ref{tab:ablation} 的 (B) 區塊納入該基線（「對稱 L2 @ $p^{*}$」一列），使兩者可分。

我們預期加權訓練仍有增益，其理由如下。Elkan 的結果成立於\emph{機率估計在整個值域上等品質}的前提；有限容量的模型並不滿足此前提，它必須決定將擬合能力配置於何處。對稱損失依樣本密度均勻配置，而決策邊界 $p^{*}=1/(1+\kappa)$ 隨 $\kappa$ 增大而深入機率分佈的尾端：於 MI300X／8B（$\kappa=6$）為 $0.14$，於平台 A／70B（$\kappa=190$）則為 $0.005$。\emph{在後者，對稱損失幾乎沒有誘因把該區域學準}——該處樣本稀少且對平均平方誤差的貢獻極低。加權訓練的作用即是將擬合能力重新配置至該區域。

此論證給出一個可證偽的預測：\textbf{(B) 區塊中「對稱 L2 @ $p^{*}$」與 (L1)/(L2) 的差距，應隨 $\kappa$ 增大而擴大}。若兩者在 3090 與 MI300X 上差距相同，則加權訓練的正當性不成立，本文將據此縮減主張至門檻校準。

\textbf{(L2) 動作成本進入排序損失的 gain 項}（原則性版本）。PARROT~\cite{parrot2020} 是我們調查中唯一刻意使用非對稱損失的工作，其 NDCG 形式明確地「重罰把機率放在即將被重用的 line 上，輕罰放在遠期重用的 line 上」，並帶來 3.5\% 的命中率提升。然而 PARROT 只有單一動作，其 gain 項 $(d(l_w)-1)$ 僅編碼\emph{重用距離}，未編碼\emph{動作成本}。我們將動作成本置入該 gain 項：
\begin{equation}
  \mathrm{DCG} = \sum_{w} \frac{\mathrm{cost}_{\text{err}}(a_w) \cdot \big(d(l_w) - 1\big)}{\log\big(\mathrm{pos}(l_w) + 1\big)}
\end{equation}

\paragraph{為何這是可辯護的新意。}
Sibyl~\cite{sibyl2022} 是唯一處理異質層級放置且考慮成本的系統，但它\emph{迴避}了此問題：其 reward 直接取自實測延遲 $1/L_t$，成本比是硬體量出來的而非猜的。這是優雅的做法，但 Sibyl 的動作\emph{可逆}（頁面可在下次存取時重新放置），而本文的 \act{Drop}$\rightarrow$重算\emph{不可逆}，因此需要 Sibyl 的驅逐懲罰項 $R_p = 0.001 \times L_e$ 推廣為 per-action 懲罰。

\paragraph{硬體條件化。}
$w(a_i)$ 由 FLOP/byte 比值決定，而該比值在 RTX 3090 為 2{,}254、MI300X 為 20{,}752（表 \ref{tab:ratio}）。\emph{同一個損失函數在兩個平台上會產生不同的權重，進而學出不同的策略}。第 \ref{sec:twoplatform} 節的硬體條件化主張，在此獲得一個具體且可量測的機制，而非僅是敘述。

此處須明確一項取捨。加權訓練使模型\emph{綁定於特定的} $\kappa$，平台變更即需重訓；上一段的兩階段做法則保持平台無關，但擬合能力的配置不隨決策邊界移動。二者並非優劣關係而是不同的操作點：前者以可移植性換取邊界附近的精度，後者反之。(B) 區塊同時量測二者，使該取捨成為\emph{可報告的結果}而非未言明的假設。

\subsection{推論開銷預算}

既有工作給出三種可行的攤銷方式：KVP~\cite{kvp2026} 於 prefill 後評分一次、decode 期間零額外 FLOP（10K 上下文下 0.71\,ms，相對於 404\,ms 的 prefill）；ForesightKV~\cite{foresightkv2026} 每 256 個新 token 重新評分一次，開銷 2.7\%；TRIM-KV~\cite{trimkv2025} 於 block 建立時評分一次，之後以指數衰減 $\beta^{t-i}$ 免費地重新排序。

我們無法直接採用 TRIM-KV 的「建立時評分一次」，因為第 \ref{sec:predictor} 節的特徵（deltas 與 EDC 計數器）\emph{在每次存取時更新}，一次性評分會立即過時。我們改採三層攤銷：（i）存取時以極低成本更新 EDC；（ii）於驅逐時機依 LRB 的做法\emph{取樣 $k=64$ 個候選評分而非排序整個 cache}，該操作在 CPU 上約需 30\,$\mu$s；（iii）粗粒度的週期性全量重新評分。

線上學習控制器能否符合真實硬體預算，Pythia~\cite{pythia2021} 提供了存在性證據：其線上 RL 硬體 prefetcher 僅佔 1.03\% 的晶片面積。

\subsection{與現有機制的關係}

\sys 不取代壓縮或稀疏注意力，而是決定\emph{何時對哪個 block 施加何種處理}。其與 vLLM 的 paged KV 佈局相容，並可與 LMCache 的分層後端並存。

% =====================================================================
\section{Evaluation Plan}
\label{sec:eval}

\textbf{本節為實驗計畫。所有結果表格均為佔位符，實驗尚未執行。}

\subsection{實驗平台}

依第 \ref{sec:twoplatform} 節的分工，實驗於兩個平台執行。

\textbf{平台 A（機制驗證與完整對照）}：NVIDIA RTX 3090（24\,GB GDDR6X，PCIe 4.0$\times$16），CUDA，vLLM。此平台承擔概念驗證、量測工具鏈建立，以及對 CUDA-only baseline（ArkVale、Quest、ShadowKV、ClusterKV、RetroInfer、ParisKV）與三個學習式 baseline 的完整比較。模型依第 \ref{sec:eval} 節的選模準則為 Qwen2.5-7B-Instruct-1M（主力）與 Llama-3.1-8B-Instruct（對照），均以 AWQ int4 權重載入。

\textbf{平台 B（規模與主要科學主張）}：AMD Instinct MI300X（gfx942，192\,GB HBM3），ROCm 7.x，vLLM ROCm wheel（Python 3.12），LMCache ROCm wheel（gfx942），真實 NVMe SSD。此平台承擔並行度／機會成本場景與能耗分析。

\paragraph{可推廣性的界線。}
我們\emph{不}由平台 A 的結果推論多租戶機會成本場景（24\,GB 無法容納多個長 session），亦\emph{不}由平台 A 提出能耗結論（消費級顯卡缺乏記憶體與計算分軌的能耗計數器）。

\subsection{壓力軸的設計}

壓力軸\emph{隨平台而不同}，這正是第 \ref{sec:twoplatform} 節雙平台設計的直接後果。

\textbf{平台 A（RTX 3090, 24\,GB）}：既有文獻的「單請求、上下文遞增」設定在此\emph{成立}。以 8-bit 權重的 Llama-3.1-8B 計，64K 上下文可置入而 128K 超出，主壓力軸為 context 8K$\rightarrow$16K$\rightarrow$32K$\rightarrow$64K$\rightarrow$（128K，預期失敗）。

\textbf{平台 B（MI300X, 192\,GB）}：同一設定在此\emph{不成立}：8B--14B 模型於 128K 僅佔 HBM 的 8--10\%（第 \ref{sec:motivation} 節），量測不到壓力。此平台改採三軸：

\begin{enumerate}\itemsep0.15em
  \item \textbf{並行度}（主軸）：1、4、8、16、32 個並行請求 @ 128K 上下文。
  \item \textbf{模型架構}：見下方選模準則。
  \item \textbf{跨 session 保留}：多輪對話與 agentic trace。
\end{enumerate}

\paragraph{選模準則。}
我們不以「涵蓋愈多模型愈好」為原則，而採一條更嚴格的規則：\emph{僅納入會使 $\kappa$ 產生可事先計算之變化的模型}。據此，選模同時受兩個硬條件約束：

\textbf{(i) 原生上下文長度必須大於評測長度。} 表 \ref{tab:kvsize} 顯示 Qwen3 全系列的 \texttt{max\_position\_embeddings} 僅 40{,}960，於 128K 評測必須施加 RoPE 外推——而外推本身降低品質，將與 KV 管理造成的品質損失混淆，使 $\epsilon$ 失去意義。此條件排除了 Qwen3 全系列。

\textbf{(ii) 必須在目標平台上留有足夠的 KV 預算。} 這排除了在平台 A 上使用 MoE：Qwen3-30B-A3B 即使以 int4 權重載入亦需約 17\,GB，僅餘 4.6\,GB 供 KV 使用。

據此，兩個平台的模型分工如下：

\begin{table}[t]
\centering\small
\setlength{\tabcolsep}{3pt}
\caption{選模結果。$^{M}$ 為 MoE。$\kappa$ 於各自平台計算。平台 A 的兩個模型 $\kappa$ 相差 2 倍，提供\emph{同一硬體上}的第二個資料點。}
\label{tab:models}
\begin{tabular}{@{}llrl@{}}
\toprule
平台 & 模型 & $\kappa$ & 角色 \\
\midrule
\multirow{2}{*}{A (3090)}
 & Qwen2.5-7B-Inst-1M & 108 & 主力（懸崖在 315K） \\
 & Llama-3.1-8B-Inst  & 54  & 對照（$\kappa$ 減半） \\
\midrule
\multirow{2}{*}{B (MI300X)}
 & Llama-4-Scout$^{M}$ & 8 & 規模（原生 10M） \\
 & Qwen3-30B-A3B$^{M}$ & 3 & MoE 對 $\kappa$ 的影響 \\
\bottomrule
\end{tabular}
\end{table}

\textbf{Qwen2.5-7B-Instruct-1M 為平台 A 的主力}，因其原生支援 1M 上下文，而其 24\,GB 上的容量懸崖落在約 315K token——\emph{遠低於模型上限，故量測到的是純粹的記憶體限制而非模型限制}。相對地，Llama-3.1-8B 在同平台的容量懸崖（約 132K）與其原生上限（131{,}072）幾乎重合，兩種限制無法分離，故僅作對照。

\paragraph{明確排除的架構。}
MLA（如 DeepSeek 系列）改變式 \eqref{eq:kvsize} 本身，其 KV 為壓縮後的潛在向量；hybrid SSM（如 Jamba）的狀態大小不隨序列成長，已由 Marconi~\cite{marconi2025} 處理。二者並非「未評估」，而是\emph{定義上不同的問題}。

\noindent 兩個平台的壓力軸不同並非實驗設計的瑕疵，而是待檢驗的假設本身：\emph{若最佳策略確實隨硬體比值改變，則產生壓力的方式本來就應該不同}。

\paragraph{上下文外推的混淆變因。}
選模準則 (i) 已排除所有需要外推的模型，故評測長度均在原生範圍內：Qwen2.5-7B-Instruct-1M 的 315K 容量懸崖遠低於其 1M 原生上限，Llama-3.1-8B-Instruct 的 131K 亦在其原生範圍內。\emph{我們因而完全避開了外推造成的品質下降這一混淆變因}，$\epsilon$ 得以純粹反映 KV 管理的影響。若日後需評測超出原生範圍的長度，必須以「full-KV + 相同外推設定」作為品質基準以隔離兩種來源。

\subsection{成本常數的量測協定}
\label{sec:costmeas}

式 \eqref{eq:cost} 的每一項都需要一組平台與模型相關的常數。這些常數\emph{每個 (平台, 模型) 組合只量一次}，之後由策略直接查用，不隨每次實驗重量。

關鍵在於\textbf{兩類成本必須分開量測}——這正是第 \ref{sec:recompute-econ} 節論證的核心，若壓成單一維度便會得到「重算最貴故永不重算」的錯誤結論：

\begin{table}[t]
\centering\small
\setlength{\tabcolsep}{3pt}
\caption{成本常數的量測協定。左欄為 block 閒置時持續付出的代價，右欄為被需要時一次性付出的代價。\act{Drop} 的價值完全來自左欄的零。}
\label{tab:costmeas}
\begin{tabular}{@{}lll@{}}
\toprule
動作 & 平時成本（閒置） & 被需要時的成本 \\
\midrule
\act{Gpu}$_{16}$ & GPU 位元組 & $\approx 0$ \\
\act{Gpu}$_{8}$  & GPU 位元組 $/2$ & 反量化 kernel \\
\act{Gpu}$_{4}$  & GPU 位元組 $/4$ & 反量化 kernel \\
\act{Cpu}$_{8}$  & host 位元組 $/2$ & PCIe + 反量化 \\
\act{Ssd}        & 磁碟位元組 & NVMe I/O $+$ PCIe \\
\act{Drop}       & \textbf{0} & 重算$(\ell)$ \\
\bottomrule
\end{tabular}
\end{table}

\paragraph{量測方式。}
平時成本以 \texttt{torch.cuda.memory\_stats()} 與 host 端計數器直接取得，單位為位元組，可由式 \eqref{eq:kvsize} 驗算。被需要時的成本則以隔離的微基準取得：對每個動作，反覆搬移或重建固定數量的 block 並取中位數延遲，重複三次確認變異小於 10\%。

\paragraph{重算成本不是常數。}
$\mathrm{recomp}(\ell)$ 隨 block 的絕對位置 $\ell$ 成長，因 attention 的二次項。我們將其量測為位置的函數而非單一純量，並以線性層的 $2N_{\text{active}}$ 項作為 $\ell \to 0$ 的下界檢核。對 MoE 模型，$N_{\text{active}}$ 為啟用專家的參數量而非總參數量。

\subsection{Baselines}

\begin{table}[t]
\centering\small
\caption{Baseline 分層與平台可用性。\textbf{平台 A（CUDA）上全部可直接執行}；平台 B 欄為移植至 ROCm 的估計工作量，經逐一檢視 repo 後得出。$^{\dagger}$ 稀疏檢索一系因 \texttt{rapidsai/raft} 無 ROCm 移植而\textbf{結構性地無法於平台 B 執行}——這是平台 A 存在的理由之一。$^{*}$ LMCache 在 ROCm 上雖有官方 wheel，但存在已知執行期問題，見第 \ref{sec:limitations} 節。ClusterKV、RetroInfer、ParisKV 同屬該系，因與上列三者機制高度重疊而未納入。}
\label{tab:baselines}
\setlength{\tabcolsep}{3pt}
\begin{tabular}{@{}llcc@{}}
\toprule
層級 & Baseline & A (CUDA) & B (ROCm) \\
\midrule
\multirow{2}{*}{錨點}
 & Full GPU（無卸載） & \checkmark & \checkmark \\
 & \textbf{Oracle（離線最優）} & \checkmark & \checkmark \\
\midrule
\multirow{4}{*}{Production}
 & vLLM \texttt{lru} & \checkmark & \checkmark \\
 & vLLM \texttt{arc} & \checkmark & \checkmark \\
 & vLLM $+$ \texttt{fs} 磁碟層 & \checkmark & \checkmark \\
 & LMCache~\cite{lmcache2025} & \checkmark & \checkmark$^{*}$ \\
\midrule
\multirow{2}{*}{學術（系統）}
 & \textbf{InfiniGen}~\cite{infinigen2024} & \checkmark & $\sim$1 週 \\
 & KIVI~\cite{kivi2024} & \checkmark & 3--7 天 \\
\midrule
\multirow{3}{*}{\textbf{學習式}}
 & \textbf{KVP}~\cite{kvp2026} & \checkmark & 3--7 天 \\
 & \textbf{ForesightKV}~\cite{foresightkv2026} & \checkmark & 3--7 天 \\
 & \textbf{LookaheadKV}~\cite{lookaheadkv2026} & \checkmark & 3--7 天 \\
\midrule
\multirow{3}{*}{稀疏檢索$^{\dagger}$}
 & ArkVale~\cite{arkvale2024} & \checkmark & $\times$ \\
 & Quest~\cite{quest2024} & \checkmark & $\times$ \\
 & ShadowKV~\cite{shadowkv2025} & \checkmark & $\times$ \\
\bottomrule
\end{tabular}
\end{table}

表 \ref{tab:baselines} 給出 baseline 集合。此處須區分兩種用途不同的對照；混淆二者會使核心主張的檢驗失去判別力。

\textbf{(i) 檢驗核心主張者為消融，而非外部系統。}本文主張既有學習式方法所用的對稱損失在異質動作成本下失準。直接比較 \sys 與 KVP 並\emph{不能}檢驗此主張：二者在儲存層數（六階對單階）、預測器族（GBDT 對強化學習）與特徵集上皆不同，量到的差異無法歸因至損失函數。具判別力的實驗是\emph{固定系統、僅替換損失函數}，即表 \ref{tab:ablation} 的 (B) 區塊，其中「對稱 L2」一列即為既有方法的損失置於本文動作空間下的行為。該實驗完全在我們控制之下，不依賴任何外部訓練產物。

\textbf{(ii) 外部學習式系統用於佐證整體競爭力。}KVP、ForesightKV 與 LookaheadKV 均已釋出程式碼，但\emph{未釋出訓練後的權重}，且其預測器為模型專屬（第 \ref{sec:limitations} 節）。我們於平台 A 重訓可得者並回報，未能重訓者引用原論文數值並明確標註未重跑。此列的完整程度不影響 (i) 的結論。

\textbf{(iii) production 策略界定實用價值。}vLLM 的 \texttt{OffloadingConnector} 內建策略\emph{僅有} \texttt{lru} 與 \texttt{arc}，其餘須經 \texttt{CachePolicy} 介面自行註冊；此二者因而構成目前 production 預設策略的完整集合。兩者皆僅依存取歷史決策：以表 \ref{tab:kappa} 的 3090／Llama-3.1-8B 為例，一個「重算需 225\,$\mu$s」的 block 與一個「自 CPU 取回需 4.2\,$\mu$s」的 block 在 LRU 眼中完全等價。本文的成本模型正是針對此空缺。

\textbf{無法於 ROCm 執行者}：ArkVale、Quest、ShadowKV、ClusterKV、RetroInfer、ParisKV。其阻擋原因並非 FlashInfer（AMD 已維護 \texttt{ROCm/flashinfer} 移植），而是 \texttt{rapidsai/raft} 無 ROCm 版本，以及 CUTLASS 的 template 層級耦合。我們將引用其論文回報數值作間接對照，並於第 \ref{sec:limitations} 節誠實說明。

\subsection{消融設計}

此為本文最重要的實驗。表 \ref{tab:ablation} 分為三個區塊，各回應一個不同的質疑。

\textbf{(A) 四項成本}回應「把四項成本組合起來是否構成貢獻」。消融\emph{以成本項為單位而非以模組為單位}：若移除任一項均導致顯著退化，其必要性即為實證結果而非主張。\textbf{我們亦承認相反的可能}：若移除某項不造成退化，正確做法是將其自系統中移除，使貢獻更為聚焦。

\textbf{(B) 損失函數}是核心主張的直接檢驗。若成本敏感損失相對於對稱 L2 無可量測優勢，則第 \ref{sec:loss} 節的整個論證不成立。此區塊不可省略。

\textbf{(C) 機制}為次要，確認品質約束與精度維度確有作用。

\paragraph{兩平台必須分開報告。}
第 \ref{sec:loss} 節論證 $w(a_i)$ 由 $\kappa$ 決定，而 $\kappa$ 在兩平台相差 32 倍（表 \ref{tab:kappa}）。(B) 區塊在兩平台上\emph{應呈現不同的退化幅度}：在 $\kappa$ 較大的 3090 上，對稱損失的代價應更明顯。\textbf{若兩平台退化幅度相同，則硬體條件化的主張被否證}——這是本文最關鍵的可證偽點。

\begin{table}[t]
\centering\small
\caption{消融矩陣。每次移除或替換一項，量測 quality--latency--memory Pareto 前緣的退化。(B) 區塊是本文核心主張的直接檢驗。$^{\dagger}$(D) 區塊另行回報\emph{熱路徑推論延遲}，因較大的預測器須在計入其自身延遲後才可比較。前三列共用攤平特徵，檢驗\emph{容量}；GRU 一列直接消費 \texttt{deltas} 序列，檢驗\emph{結構}。任一列在計入自身延遲後仍佔優，第 \ref{sec:predictor} 節的對應設計選擇即被推翻。}
\label{tab:ablation}
\setlength{\tabcolsep}{3.5pt}
\begin{tabular}{llcc}
\toprule
 & & \multicolumn{2}{c}{Pareto 退化} \\
\cmidrule(l){3-4}
變體 & 移除／替換 & 3090 & MI300X \\
\midrule
\multicolumn{4}{l}{\textit{\textbf{(A) 四項成本}}} \\
\sys（完整）        & ---                        & \placeholder{基準} & \placeholder{基準} \\
$-$ 未來效用        & $p_i \rightarrow$ recency  & \placeholder{TBD} & \placeholder{TBD} \\
$-$ 傳輸成本        & $\lambda_x = 0$            & \placeholder{TBD} & \placeholder{TBD} \\
$-$ 重算成本        & 移除 \act{Drop}            & \placeholder{TBD} & \placeholder{TBD} \\
$-$ 記憶體成本      & $\lambda_m = 0$            & \placeholder{TBD} & \placeholder{TBD} \\
\midrule
\multicolumn{4}{l}{\textit{\textbf{(B) 損失函數}（核心主張）}} \\
對稱 L2 @ $0.5$     & $w \equiv 1$，門檻 $0.5$   & \placeholder{TBD} & \placeholder{TBD} \\
\textbf{對稱 L2 @ $p^{*}$} & \textbf{$w \equiv 1$，門檻式 \eqref{eq:threshold}} & \placeholder{TBD} & \placeholder{TBD} \\
(L1) 加權 + 校準    & 式 \eqref{eq:threshold}    & \placeholder{TBD} & \placeholder{TBD} \\
(L2) 成本進 NDCG    & ---                        & \placeholder{TBD} & \placeholder{TBD} \\
\midrule
\multicolumn{4}{l}{\textit{\textbf{(C) 機制與特徵}}} \\
$-$ 品質約束        & $\epsilon \rightarrow$ 固定預算 & \placeholder{TBD} & \placeholder{TBD} \\
$-$ 壓縮            & 移除精度維度               & \placeholder{TBD} & \placeholder{TBD} \\
$-$ 模型內部特徵    & 僅存取歷史                 & \placeholder{TBD} & \placeholder{TBD} \\
$-$ query 側訊號    & 移除 \texttt{attn\_mass}   & \placeholder{TBD} & \placeholder{TBD} \\
\midrule
\multicolumn{4}{l}{\textit{\textbf{(D) 預測器容量與結構}}$^{\dagger}$} \\
GBDT（本文）        & ---                        & \placeholder{基準} & \placeholder{基準} \\
MLP $2\times64$     & 容量：同特徵、同損失       & \placeholder{TBD} & \placeholder{TBD} \\
MLP $4\times256$    & 容量：同特徵、同損失       & \placeholder{TBD} & \placeholder{TBD} \\
\textbf{GRU}（序列） & \textbf{結構}：直接吃 \texttt{deltas} & \placeholder{TBD} & \placeholder{TBD} \\
\bottomrule
\end{tabular}
\end{table}



\subsection{權重精度的敏感度分析}

權重精度不是本文的決策變數，但它改變 KV 可用的預算，因而間接影響最佳策略。我們將其列為敏感度分析而非主要結果：

\begin{center}\footnotesize
\setlength{\tabcolsep}{3pt}
\begin{tabular}{@{}lll@{}}
\toprule
平台 & 權重 & 目的 \\
\midrule
A (3090) & AWQ-INT4 & 主要設定（記憶體受限） \\
A (3090) & BF16 & 對照：量化對 KV 預算的影響 \\
B (MI300X) & BF16 & 主要設定 \\
B (MI300X) & FP8 & 對照：CDNA3 之後的常見配置 \\
\bottomrule
\end{tabular}
\end{center}

\textbf{待檢驗的問題}：在固定的硬體與模型下，最佳 KV 策略是否隨權重精度改變？我們預期\emph{會}——因為權重精度決定 KV 預算，而預算改變會移動式 \eqref{eq:opt} 的容量約束。但這是\emph{間接}影響，與第 \ref{sec:twoplatform} 節的 $\kappa$ 軸（直接改變成本比）性質不同。

\paragraph{一個須澄清的區分：FP8 運算與 FP8 儲存。}
Ampere（sm\_86）確無原生 FP8 \emph{tensor core 運算}，但本文動作空間中的 \act{GPU-FP8} 是\emph{儲存}狀態——KV 以 FP8 存放、讀取時反量化——與 tensor core 無關。我們在平台 A 上實測 \texttt{-{}-kv-cache-dtype fp8} 可正常運作：Llama-3.1-8B 的 KV 容量由 41{,}648 增至 83{,}312 token、Qwen2.5-7B-1M 由 106{,}512 增至 213{,}040 token，二者皆為\emph{恰好兩倍}且位元組佔用不變。

更進一步，vLLM 亦提供 \texttt{int8\_per\_token\_head} 與 \texttt{int4\_per\_token\_head}，故\textbf{動作空間中位於 GPU 上的四個精度階，在平台 A 上皆為同一旗標的不同取值}，無須額外實作。以每 GiB 可容 token 數正規化後（消除 KV pool 大小的 run-to-run 抖動，全距降至 0.04\%）實測得：

\begin{center}
\small
\begin{tabular}{lrrrr}
\toprule
KV dtype & KiB/tok & 理想 & 額外 & vs BF16 \\
\midrule
\texttt{auto} (BF16) & 128.11 & 128 & $+0.11$ & 1.00$\times$ \\
\texttt{fp8} & 64.04 & 64 & $+0.04$ & \textbf{2.00}$\times$ \\
\texttt{int8\_per\_tok\_head} & 66.04 & 64 & $\mathbf{+2.04}$ & 1.94$\times$ \\
\texttt{int4\_per\_tok\_head} & 34.02 & 32 & $\mathbf{+2.02}$ & \textbf{3.77}$\times$ \\
\bottomrule
\end{tabular}
\end{center}

\texttt{fp8} 為純格式轉換故恰為 $2\times$；兩個 per-token-head 量化階則各固定多出 2\,KiB/token。此開銷可由架構直接算出：Llama-3.1-8B 有 32 層、8 個 KV head，K 與 V 各一份，故每 token 有 $2\times32\times8=512$ 個 scale，以 FP32 儲存即 $512\times4=2{,}048$\,bytes $=2$\,KiB，與實測相符（int8 為 4.1\,bytes/scale、int4 為 4.0）。\textbf{故精度階梯的收益並非 $2\times$/$4\times$，而是 $2.00\times$/$1.94\times$/$3.77\times$}；式 \eqref{eq:kvsize} 若僅以每元素位元組數計算，將系統性高估低精度階的收益。量化中繼資料是與序列長度成正比的固定成本，必須進入成本模型。

\subsection{Oracle 上界與 go/no-go 判準}

我們建構一個離線最優 oracle：給定完整的未來查詢序列，以整數規劃求解式 \eqref{eq:opt}，得到 $p_i \in \{0,1\}$ 為已知時的最佳放置。此為經典快取文獻中 Belady~\cite{belady1966} 分析的直接延伸，並與 Sibyl~\cite{sibyl2022} 報告「達到全知 oracle 的 80\%」的評估傳統一致。

Fang 等人~\cite{hetmem2025} 已為兩層（HBM 與 off-package DRAM）的無損放置推導上界。我們的 oracle 在三點上延伸之：納入 SSD 層、納入 \act{Drop}$\rightarrow$重算動作、納入品質約束。

\paragraph{停止條件。}
若 oracle 相對於最佳簡單策略的改善幅度小於 5\%，我們將終止此研究方向並報告此負面結果。此判準於實驗計畫的第 5--6 週執行，早於任何策略實作。

\subsection{品質約束的兩類動作}
\label{sec:quality-two-classes}

式 \eqref{eq:opt} 的 $\epsilon$ 約束對動作空間中的兩類動作意義不同，量法亦須不同。

\paragraph{無損動作：$\epsilon=0$ 是可驗證的斷言，不是假設。}
\act{CPU} 與 \act{SSD} 僅搬移位元組，理應不改變輸出。此為\emph{正確性}斷言而非品質取捨：若輸出有異，即為實作錯誤。我們以逐字元比對驗證之——以 64-shot GSM8K（共用前綴 11{,}029 token）於平台 A 對 60 題取輸出的 SHA-1，與 \act{GPU-BF16} 基準比對：

\begin{center}
\small
\begin{tabular}{lrr}
\toprule
設定 & 正確率 & 逐字元相同 \\
\midrule
\texttt{full\_gpu}（基準） & 76.67\% & --- \\
\texttt{cpu\_lru} & 76.67\% & \textbf{60/60} \\
\texttt{tier\_fs}（CPU$+$磁碟） & 76.67\% & \textbf{60/60} \\
\bottomrule
\end{tabular}
\end{center}

\emph{位置分層的三階在 $\epsilon$ 約束下等價}；$\epsilon$ 的實質內容全部落在精度分層上。

\paragraph{有損動作：目前的樣本數不足以分辨各精度階。}
我們以同一工作負載量測精度階梯（$n=120$，貪婪解碼，固定 seed）：

\begin{center}
\small
\begin{tabular}{lrrr}
\toprule
KV dtype & 正確率 & 相對 BF16 & $z$ \\
\midrule
\texttt{auto} (BF16) & 79.17\% & --- & --- \\
\texttt{fp8} & 73.33\% & $-5.83$\,pt & 1.11 \\
\texttt{int8\_per\_tok\_head} & 80.00\% & $+0.83$\,pt & 0.16 \\
\texttt{int4\_per\_tok\_head} & 75.83\% & $-3.33$\,pt & 0.64 \\
\bottomrule
\end{tabular}
\end{center}

\textbf{三個差值皆不顯著。} $n=120$、$p\approx0.79$ 時單組標準誤為 3.71\,pt，兩組差值的標準誤為 5.24\,pt，95\% 信賴區間為 $\pm10.3$\,pt；三個 $z$ 值均遠低於 1.96。差值亦非單調（\texttt{int8} 高於 BF16，而 \texttt{fp8} 低於 \texttt{int4}），此非單調性本身即為樣本量不足的徵候。

\emph{我們因此不宣稱任何精度階之間存在品質差異。} 欲以 80\% 檢定力分辨 5\,pt 的差距（$\alpha=0.05$）需每組約 1{,}055 題；GSM8K 測試集有 1{,}319 題，故此量測可行，將於定稿版本補足。在此之前，表 \ref{tab:results} 中 $\epsilon$ 相關的欄位一律標記為 \texttt{NOT\_MEASURED}。

\subsection{評估指標與工作負載}

系統指標：TTFT、TPOT、端到端延遲、GPU/CPU/SSD 佔用、PCIe traffic、throughput、服務層級分佈（見下）、重算成本。

\paragraph{ROCm 量測方法。}
KV 搬移經由 SDMA copy engine 而非 AQL compute queue，\emph{對 \texttt{rocprof-compute} 的 TCC/L2 counter 不可見}。我們將使用 \texttt{rocprofv3 --memory-copy-trace} 取得逐筆傳輸的位元組數、方向與時間戳，並以 \texttt{rocprof-compute} 的 gfx942 "Remote Read Traffic" 指標取得 kernel 側視角。

\paragraph{能耗。}
MI300X 的 \texttt{amdsmi} 提供 per-rail 能耗累加器，其中 \texttt{hbm\_energy\_acc} \emph{將 HBM 能耗與計算能耗分離}。對一個探討「保留於記憶體或重新計算」的研究問題而言，這是直接對應的量測能力，且在 NVIDIA 平台上不可得。

\paragraph{品質指標。}
Yandex 的研究~\cite{yakv2026} 顯示 ShadowKV 等方法在 needle-in-a-haystack 上表現良好卻在\emph{多事實抽取}任務上失效。因此我們將 needle-in-a-haystack 降級為 sanity check，改以多事實抽取作為主要品質判準，另納入 LongBench、RULER 與長文本摘要。

\paragraph{品質如何化約為單一數字。}
表 \ref{tab:results} 的 Quality 欄須為單一純量，其定義如下。設 $\mathcal{T}$ 為上述任務集合，$Q_t(\pi)$ 為策略 $\pi$ 於任務 $t$ 的原始分數，並以 Full GPU 為參照定義相對保留率 $r_t(\pi) = Q_t(\pi)\,/\,Q_t(\textsf{FullGPU})$。我們回報
\begin{equation}
Q(\pi) \;=\; \min_{t \in \mathcal{T}} \; r_t(\pi),
\label{eq:quality}
\end{equation}
即\emph{最差任務}的保留率，而非跨任務平均。理由與上一段相同：平均會使單一任務的崩潰被其餘任務的分數掩蓋，而 \cite{yakv2026} 所報告的正是此類崩潰。品質約束因而具體化為 $Q(\pi) \ge 1-\epsilon$。逐任務的 $r_t$ 完整列於附錄，避免式 \eqref{eq:quality} 的化約隱藏任務間的差異。

\paragraph{以服務層級分佈取代單一命中率。}
「KV hit rate」對單層快取有明確意義，但本文的動作空間有六階，「命中」並非二元事件：自 CPU 取回與自 SSD 取回的成本相差約一個數量級，與重算則相差 6--190 倍（表 \ref{tab:kappa}）。將這些情形壓成單一比率會抹除本文所要量測的成本結構。我們因此回報\emph{服務層級分佈}，即每次 KV 存取由 $\mathcal{A}$ 中哪一階滿足的經驗分佈，並另外給出兩個純量：(i) \textbf{GPU 常駐率}，存取時該 block 已位於 GPU 且為 BF16 的比例（對應傳統命中率）；(ii) \textbf{重算率}，由 \act{Drop} 觸發重算的存取比例。後者為本文特有，直接對應第 \ref{sec:recompute-econ} 節的成本模型，亦是表 \ref{tab:ablation} 中「移除 \act{Drop}」一列的解讀依據。

\paragraph{預測器層級的診斷：為何不回報分類準確度。}
我們\emph{刻意不}以預測器的分類準確度作為評估指標，並在此說明理由，以免被誤讀為疏漏。成本敏感訓練的作用即是將決策邊界自 $0.5$ 移至 $p^{*}=1/(1+\kappa)$，其必然後果是模型在\emph{廉價方向}上多犯錯以換取在\emph{昂貴方向}上少犯錯。以未加權的準確度衡量，成本敏感模型\emph{應當}劣於對稱模型；若其未劣化，反而意味著加權未生效。準確度因而與本文的目標函數方向相反，不構成有效證據。我們改以兩個量診斷預測器：

\emph{(i) 校準品質。}式 \eqref{eq:threshold} 的最佳門檻僅在 $\hat{p}_i$ 為\emph{校準過的}機率時成立：若模型輸出 $\hat p = 0.1$ 的樣本中實際重用比例並非約 $10\%$，該門檻即失去意義，第 \ref{sec:loss} 節的整個論證隨之失效。我們回報 isotonic regression 校準前後的可靠度圖與 expected calibration error。此量同時是第 \ref{sec:loss} 節所揭露之 L1 缺陷（加權迴歸改變條件均值而非決策門檻）是否被校準步驟補救的\emph{唯一直接證據}；缺少此圖，該缺陷等同未被處理。

\emph{(ii) 成本加權錯誤率。}$\sum_i c\!\left(a_i^{\text{pred}},\, a_i^{\star}\right)$，其中 $a_i^{\star}$ 為 oracle 動作、$c(\cdot,\cdot)$ 為式 \eqref{eq:cost} 的成本函數。此量計的是「錯掉多少微秒」而非「錯了幾次」，是唯一能在預測器層級公平比較兩種損失函數的指標，亦與系統層級的端到端延遲直接可對照。

\paragraph{跨工作負載泛化。}
承第 \ref{sec:predictor} 節的分布飄移討論，我們額外回報一組交叉實驗：於單一工作負載訓練並於其餘工作負載測試，回報上述兩個診斷量與 $Q(\pi)$ 的退化幅度。此結果同時界定週期性重訓的必要性——若跨工作負載退化可忽略，重訓機制即為不必要的複雜度，我們將據以移除。

\subsection{已完成的量測推翻了本文的四項前提}
\label{sec:measured}

我們在平台 A 上完成成本常數與容量的量測。四項結果與本文原先的敘述相牴觸，其中兩項擴大了可驗證的範圍，兩項要求修正成本模型。表 \ref{tab:measured} 列出每項量測、其推翻的敘述、以及對論文的後果；量測協定與逐設定的完整數據見附錄 \ref{app:measurements}。

\begin{table*}[t]
\centering\small
\caption{平台 A（RTX 3090, sm\_86）的已完成量測。每列給出一項與本文原敘述相牴觸的結果。}
\label{tab:measured}
\begin{tabular}{p{0.20\textwidth}p{0.30\textwidth}p{0.42\textwidth}}
\toprule
\textbf{量測} & \textbf{原敘述 $\rightarrow$ 實測} & \textbf{對本文的後果} \\
\midrule
FP8 KV 於 sm\_86 &
「Ampere 無原生 FP8，平台 A 無此設定」$\rightarrow$ 可用，容量恰增 $2.00\times$（41{,}648$\rightarrow$83{,}312 token） &
混淆了 FP8 \emph{運算}與 FP8 \emph{儲存}。動作空間中位於 GPU 的四個精度階皆為同一旗標的取值，單一平台即可驗證五階 \\
\addlinespace
精度階梯的實際收益 &
$2\times$/$4\times$ $\rightarrow$ $2.00\times$/$1.94\times$/$3.77\times$ &
per-token-head 量化固定多帶 2\,KiB/token 的 scale（$2\times32\times8$ 個 FP32），式 \eqref{eq:kvsize} 僅計每元素位元組數會高估低精度階 \\
\addlinespace
$C_{\text{recompute}}$ 對位置的形式 &
「attention 的二次項」$\rightarrow$ 對位置線性，式 \eqref{eq:recompute-position}，擬合偏差 $<1.6\%$ &
整段 prefill 為 $O(L^2)$，但固定大小 block 在位置 $P$ 重算為 $O(C\cdot P)$。原敘述混用了兩個量 \\
\addlinespace
\act{SSD} 與 \act{Drop} 的優劣次序 &
「兩者同數量級」$\rightarrow$ 成立，且\emph{次序隨 block 的絕對位置反轉}，交叉點在 7{,}277\,token &
\act{Drop} 的成本隨位置線性成長而 \act{SSD} 為常數，故位置 $<7{,}277$ 時 \act{Drop} 較廉、其後 \act{SSD} 較廉。於 128K 上下文，僅 5.6\% 的 block 落在 \act{Drop} 一側。\textbf{動作空間的有效子集隨上下文長度而變}——此即本文主張的一個實例 \\
\bottomrule
\end{tabular}
\end{table*}

\paragraph{Takeaway.}
四項結果中，兩項擴大了平台 A 的可驗證範圍（FP8 儲存可用、四個精度階同旗標可切），兩項要求修正成本模型（量化中繼資料、重算的位置依賴形式）。第四項的性質不同：它不是修正，而是本文核心主張的直接實例——\act{SSD} 在平台 A 被 \act{Drop} 支配，故六階動作空間並非普遍適用。論文若宣稱六階皆為必要，此結果即為反例；改為「動作空間的\emph{有效}子集隨平台而定」則與證據一致，且更強。

\subsection{權重精度決定平台 A 能否進入長上下文區間}
\label{sec:awq-unlock}

BF16 權重在 24\,GB 卡上佔 15\,GB，KV 預算僅餘 5.9\,GB，容量上限 41{,}648 token。此配置\emph{無法}進入本文關心的 128K 區間。AWQ-INT4 權重佔 4.7\,GB，KV 預算增至約 15\,GB：

\begin{center}
\footnotesize
\begin{tabular}{@{}lrrr@{}}
\toprule
設定 & KV 容量 & 定址上限 & 有效上限 \\
\midrule
Llama-8B BF16 權重 & 41{,}648 & 131{,}072 & 41{,}648 \\
Llama-8B AWQ 權重 & 120{,}320 & 131{,}072 & 120{,}320 \\
Llama-8B AWQ $+$ FP8-KV & 240{,}656 & 131{,}072 & \textbf{131{,}072}$^\dagger$ \\
Qwen-7B-1M AWQ 權重 & 273{,}872 & 262{,}144 & \textbf{262{,}144}$^\dagger$ \\
Qwen-7B-1M AWQ $+$ FP8-KV & \textbf{547{,}744} & 262{,}144 & \textbf{262{,}144}$^\dagger$ \\
DeepSeek-V2-Lite MLA AWQ & \textbf{399{,}376} & 163{,}840 & \textbf{163{,}840}$^\dagger$ \\
\bottomrule
\end{tabular}
\end{center}

$^\dagger$ 記憶體容量超出模型的可定址長度，瓶頸由硬體轉為架構。

\paragraph{Takeaway.}
在正確的權重配置下，24\,GB 消費級 GPU 的 KV 記憶體\emph{不是}長上下文的瓶頸；模型的可定址長度才是。六個設定中有四個的 KV 容量超出模型可定址長度，其中 Qwen-7B-1M 於 AWQ 權重加 FP8-KV 下達 547{,}744 token，為其可定址上限的 $2.1$ 倍。此點對本文的雙平台設計有兩個後果。其一，平台 A 足以驗證 128K 區間的容量與成本主張，無須更大的加速器。其二，MLA 架構將 KV 壓至 30.4\,KiB/token（GQA 的 $1/4.2$），使單卡容量達 399{,}376 token——\emph{單請求的容量壓力於此消失}，壓力轉為 PCIe 頻寬與多租戶下的機會成本，即第 \ref{sec:twoplatform} 節所述平台 B 的情境。此結果將 MLA 納入本文的 $\kappa$ 軸，且在單一平台上提供 $4.2\times$ 的跨度。

\subsection{預期結果（佔位）}

\begin{table}[t]
\centering\small
\caption{\placeholder{主結果表（佔位符，實驗尚未執行）}}
\label{tab:results}
\begin{tabular}{lccc}
\toprule
方法 & Quality & TTFT & GPU mem \\
\midrule
Full GPU        & \placeholder{---} & \placeholder{---} & \placeholder{---} \\
vLLM \texttt{lru} & \placeholder{---} & \placeholder{---} & \placeholder{---} \\
vLLM \texttt{arc} & \placeholder{---} & \placeholder{---} & \placeholder{---} \\
LMCache         & \placeholder{---} & \placeholder{---} & \placeholder{---} \\
InfiniGen       & \placeholder{---} & \placeholder{---} & \placeholder{---} \\
\sys            & \placeholder{---} & \placeholder{---} & \placeholder{---} \\
\midrule
\textit{Oracle} & \placeholder{---} & \placeholder{---} & \placeholder{---} \\
\bottomrule
\end{tabular}
\end{table}

% =====================================================================
\section{Discussion and Limitations}
\label{sec:limitations}

\paragraph{本稿無實驗結果。}
這是本稿最主要的限制。第 \ref{sec:motivation} 節的分析與第 \ref{sec:related} 節的比對已完成，但第 \ref{sec:eval} 節的計畫尚未執行。在取得 oracle 結果前，我們無法斷言此問題存在足夠的最佳化空間。

\paragraph{$\epsilon$ 是估計值而非保證。}
品質約束以 profiling 建立的代理訊號執行，$\Delta\mathcal{Q} \le \epsilon$ 是經驗性的估計而非形式保證。這是所有品質約束工作的共同限制。將其提升為執行期可驗證的界限（例如以 per-head 的注意力分佈失真上界）是自然的後續方向，但其理論難度高。

\paragraph{重算可能不划算。}
第 \ref{sec:recompute-econ} 節的分析顯示重算的單次成本比傳輸高 6--190 倍。若 $p_i$ 的預測精度不足，\act{Drop} 動作可能持續造成淨損失。消融實驗（表 \ref{tab:ablation}）將直接檢驗此點；若 \act{Drop} 不帶來收益，我們將自系統中移除之。

\paragraph{K 與 V 未分離處理。}
式 \eqref{eq:kvsize} 前導的 2 對應 key 與 value，本文\emph{將二者綁定於同一動作}：一個 block 的 K 與 V 一同降級、一同量化、一同丟棄。這是刻意的範圍限制，而文獻指出該粒度並非最佳。KIVI~\cite{kivi2024} 顯示二者的元素分佈不同，key cache 宜沿 channel 維度量化而 value cache 宜沿 token 維度量化；Hariri 等人~\cite{kvadaquant2025} 進一步指出在固定記憶體預算下，將位元優先配置給 key（如 4-bit key 搭配 2-bit value）較均勻配置嚴格更優；KVTuner~\cite{kvtuner2025} 則於逐層粒度搜尋混合精度，在 Llama-3.1-8B-Instruct 上達成 3.25 位元的近無損配置。

因此本文階梯上的 \act{Gpu}$_{8}$ 與 \act{Gpu}$_{4}$ 應理解為 K/V 共用精度的\emph{保守設定}，其品質退化為 K/V 分離配置的上界。將動作空間沿 K/V 拆為兩軸是自然的延伸，但會使 $|\mathcal{A}|$ 自 6 增至 36，並使標籤生成與 oracle 求解的成本相應上升。本文不涵蓋此延伸；它與本文的核心主張正交——無論動作以何種粒度定義，「動作成本異質時損失函數須編碼該成本」的論點皆成立，且動作粒度變細只會使成本的離散程度更高，因而強化而非削弱該論點。

\paragraph{ROCm 生態的成熟度。}
LMCache 無自動化的 on-device ROCm CI；其預設 MP 模式在 ROCm 上存在已知的死鎖問題；vLLM 在 ROCm 上的 KV 傳輸錯誤目前會導致 worker 終止而非降級（相關修正 PR 尚未合併）。此外，啟用 KV connector 會使 vLLM 自候選清單移除 \texttt{ROCM\_ATTN} 後端，我們將記錄每次實驗實際使用的 attention 後端。

\paragraph{與 AMD 生產堆疊的關係。}
AMD 的 \texttt{ROCm/mori} 中的 MORI-UMBP 提供「具分層儲存與分散式鍵值存取的統一記憶體與頻寬池」，已整合入 SGLang HiCache 並具備 SSD 層。我們不宣稱 AMD 平台上缺乏分層 KV 基礎設施；本文的貢獻在於其\emph{放置策略}，而非分層機制本身。

\paragraph{平台 A（RTX 3090）的限制。}
24\,GB 的容量意味著平台 A \emph{無法}執行單請求 128K 上下文（第 \ref{sec:twoplatform} 節），亦無法承載多租戶場景，本文不由平台 A 推論機會成本相關結論。此外消費級顯卡缺乏記憶體與計算分軌的能耗計數器，能耗分析僅於平台 B 進行。

\paragraph{學習式對照組未釋出權重。}
KVP~\cite{kvp2026}、ForesightKV~\cite{foresightkv2026} 與 LookaheadKV~\cite{lookaheadkv2026} 三者\emph{均未釋出訓練後的權重}，且其預測器皆為模型專屬——KVP 需為 Qwen2.5-7B-Instruct 的 112 個 layer-head 各訓練一個 agent（其 README 載明使用 8 GPU DDP），ForesightKV 需至少兩張 GPU 分置訓練模型與參考模型，LookaheadKV 僅提供「minimal training example」。重訓因而是\emph{每個主模型各一次的完整訓練}，且主模型變更時須全部重做。

我們因此\emph{不}以外部學習式系統作為核心主張的檢驗手段。本文的主張是「對稱損失在異質動作成本下失準」，而檢驗該主張的正確方式是\emph{在同一系統內替換損失函數}（表 \ref{tab:ablation} 的 (B) 區塊），此設計消除了系統間的無關差異。外部學習式 baseline 僅用於佐證整體競爭力，其數值取自原論文並明確標註未重跑。

\paragraph{無法重跑的對照組。}
ArkVale、Quest、ShadowKV、ClusterKV、RetroInfer 與 ParisKV 因 \texttt{rapidsai/raft} 無 ROCm 移植而無法於平台 B 執行；我們於平台 A 執行之，但\emph{無法提供這些方法在 MI300X 上的數值}。

\paragraph{MI300A 的架構差異。}
若目標平台為 MI300A（APU）而非 MI300X，本文的階梯需重新設計。MI300A 的 CPU 與 GPU 共用同一塊 128\,GB HBM3，單一 OS 管理的實體位址空間，無 DIMM 亦無主機 DDR~\cite{cdna3,upm2025}。「將 KV 卸載至 CPU 記憶體」在該架構上僅為同一記憶體內的頁表變更，不產生容量收益；階梯將塌縮為 HBM $\rightarrow$ 對等 APU HBM $\rightarrow$ NVMe。

% =====================================================================
\section{Conclusion}

本文的核心論點是：\emph{當動作成本相差 6--190 倍、且該比值本身隨硬體變動 32 倍時，以對稱損失訓練的效用預測器會系統性地錯置決策邊界。}我們將 KV 放置形式化為受品質約束的序列決策問題，其動作空間統一了位置分層與精度分層，並將「丟棄後重算」視為階梯上一個具明確標價的層級。

平台 A 的量測支持此論點，並在四處修正了本文的敘述（第 \ref{sec:measured} 節）。其中兩處擴大了可驗證範圍：FP8 為\emph{儲存}狀態而非運算狀態，故 sm\_86 可量測，且動作空間中位於 GPU 的四個精度階皆為同一旗標的取值。另兩處要求修正成本模型：per-token-head 量化固定多帶 2\,KiB/token 的中繼資料，使精度階梯的收益為 $2.00\times$/$1.94\times$/$3.77\times$ 而非 $2\times$/$4\times$；而固定大小 block 的重算成本對其絕對位置為\emph{線性}，非二次。

第四處修正的性質不同。\act{SSD} 的取回成本為常數（5.54\,ms/block），而 \act{Drop} 的成本隨 block 的絕對位置線性成長（式 \eqref{eq:recompute-position}），二者於位置 7{,}277\,token 處交叉：其前 \act{Drop} 較廉，其後 \act{SSD} 較廉。於 128K 上下文，僅 5.6\% 的 block 落在 \act{Drop} 一側；於 512K 則為 1.4\%。這不是本文主張的反例，而是其直接實例：\emph{最佳動作空間的有效子集隨上下文長度與平台共同決定}，而二者皆可事先算出。六階動作空間應理解為候選集合。

量測亦劃清了品質約束的作用範圍。\act{CPU} 與 \act{SSD} 的輸出與 \act{GPU-BF16} 逐字元相同（60/60），故位置分層的三階在 $\epsilon$ 約束下等價；$\epsilon$ 的實質內容全部落在精度分層上。我們目前的樣本量（$n=120$）不足以分辨各精度階之間的品質差異，故不作此宣稱。

尚未完成的部分決定了本文的下一步。策略的實作與 Oracle 的 go/no-go 判定尚待乾淨環境下的量測；本文的 $\kappa$ 跨硬體主張需要平台 B 的數據；而多租戶下的機會成本情境在 24\,GB 卡上無法驗證——單卡放不下兩個 128K session。第 \ref{sec:eval} 節給出的判準與第 \ref{app:oracle} 節給出的求解方法效力範圍共同約束了這些後續工作可以宣稱什麼。

% =====================================================================
\section*{Statements}

\paragraph{Data Availability.}
本稿為初稿，尚無實驗資料。第 \ref{sec:motivation} 節的 KV footprint 計算工具將隨程式碼一併釋出。所引用模型組態均取自公開的 HuggingFace \texttt{config.json}。

\paragraph{Ethics.}
本研究不涉及人類受試者或個人資料。工作負載將使用公開 benchmark 與已匿名化的公開對話 trace。

\paragraph{Author Contributions (CRediT).}
\textit{（投稿時填入）} Conceptualization; Methodology; Software; Investigation; Writing --- original draft.

\paragraph{Conflict of Interest.}
\textit{（投稿時填入）}

\paragraph{Funding.}
\textit{（投稿時填入）}

\paragraph{AI Usage Disclosure.}
本稿的文獻調查、比對分析與初稿撰寫過程使用了 AI 輔助研究工具。所有引用文獻均經作者查證其存在性與內容；所有量化分析均可由公開規格與模型組態重現。AI 工具未用於產生實驗數據。

% =====================================================================
\appendix

\section{平台 A 的量測明細}
\label{app:measurements}

本附錄給出第 \ref{sec:measured} 節所引數字的量測協定與逐設定數據。所有量測於單張 RTX 3090（sm\_86, driver 550.163.01）上以 vLLM 0.28.0+cu129 執行，每列數據可由 \texttt{run\_id} 追溯至指令與輸出檔。

\subsection{量測衛生}
\label{app:hygiene}

該機器由二十餘位使用者共用，此事實決定了三項協定。

\paragraph{每卡獨佔不足以保證乾淨。}
GPU 的 SM 為各卡獨佔，但 PCIe、host RAM 頻寬與 \texttt{/dev/shm} 為全機共用。他人於其他卡上的負載不會出現在本卡的 compute-apps 中，卻會拖慢搬移量測。我們逐格比較序列執行與五卡並行執行的 80 個量測點：不經 PCIe 搬移 KV 的 \texttt{full\_gpu} 差異在 $\pm2\%$ 內，而經 PCIe 的 warm TTFT 差異達 $-26\%$ 至 $-52\%$。故每列數據記錄整機爭用等級，且僅序列模式的數據進入論文。

\paragraph{KV pool 大小在同一設定下並非常數。}
同一模型、同一卡、同一 \texttt{gpu\_memory\_utilization}，\texttt{GPU KV cache size} 在 5.09 與 5.88\,GiB 兩值間擺動（原始 token 數全距 13.5\%）。以每 GiB 可容 token 數正規化後全距降至 0.04\%，故容量的平時成本一律以 KiB/token 報告，而非原始 token 數。前一行程結束至驅動釋放記憶體之間存在延遲，我們要求連續三次取樣皆達標才啟動下一次量測。

\paragraph{分層量測須驗證資料確實落到目標層。}
量測 \act{SSD} 時若 CPU 階容量大於工作集，資料不會 cascade 至磁碟。我們的首次量測即因此失效：CPU 階 24\,GiB、工作集 8\,GiB，vLLM 的 \texttt{chunk\_queries} 顯示磁碟階僅被查詢 2 次而 CPU 階 2{,}048 次，所得「\act{SSD} 與 \act{CPU} 僅差 1.2\%」實為 CPU 階的數字。將 CPU 階縮至 1\,GiB 後磁碟階被查詢 2{,}048 次，實際讀取 3\,GiB。每列數據記錄兩層的 \texttt{chunk\_queries}。

\subsection{成本常數}
\label{app:cost}

\paragraph{平時成本。}
以每 GiB 可容 token 數量測，每設定重複三次取中位數（全距 $\le0.04\%$）：BF16 128.11、\texttt{fp8} 64.04、\texttt{int8\_per\_tok\_head} 66.04、\texttt{int4\_per\_tok\_head} 34.02\,KiB/token。後兩者各多出約 2.03\,KiB/token，與 $2\times32\times8=512$ 個 FP32 scale 的 2{,}048\,bytes 相符（實測每 scale 4.0--4.1\,bytes）。

\paragraph{被需要時的成本。}
以 16K 上下文、四個共用前綴的兩輪工作負載量測 warm TTFT，並扣除 GPU 命中的基準（134.4\,ms）：\act{CPU} 736.9\,ms、\act{SSD} 5{,}803.3\,ms、\act{Drop} 5{,}467.5\,ms。換算每 block（16 token）分別為 0.588、5.536、4.008\,ms。

\paragraph{磁碟階的成本並非由 I/O 主導。}
我們以 \texttt{O\_DIRECT} 循序存取量測兩顆裝置，各重複三次取中位數。SATA QLC（Samsung 870 QVO，\texttt{/ssd7}）讀 522\,MiB/s，NVMe（Crucial P3，\texttt{/}）讀 2{,}085\,MiB/s，相差 $4.0\times$；端到端 warm TTFT 卻是 5{,}803 與 6{,}600\,ms，頻寬較高者反而較慢。由端到端延遲反推的 \act{SSD} 取回成本 5.536\,ms/block 換算為 344\,MiB/s，僅達裸讀的 66\%，其餘為分層機制的查表、cascade 排程與 promotion 開銷。成本模型因此拆為 $C_{\text{ssd}}=C_{\text{lookup}}+C_{\text{sched}}+C_{\text{io}}(\text{device})$，僅末項與裝置有關。兩條互相獨立的路徑（端到端 TTFT 與原始 I/O）在 66\% 這個比例上互相印證。

\paragraph{QLC 的持續寫入只有短測的 37\%。}
1\,GiB 的短測給出 492\,MiB/s，16\,GiB 的長測只剩 181\,MiB/s；短測整個落在 QLC 的 SLC 寫入快取內。KV 階持續寫入，可行性判定因此須採用 181\,MiB/s。同一台機器上的 NVMe 寫入為 2{,}512\,MiB/s，兩者相差 $13.9\times$。\act{SSD} 這個動作的成本因此在單一台機器內就相差一個數量級，而表 \ref{tab:kappa} 的 $\kappa$ 僅涵蓋加速器與主機之間的傳輸，未涵蓋磁碟階。

\subsection{Oracle 的求解方法與其效力範圍}
\label{app:oracle}

我們以 trace 驅動的模擬求解上界：GPU 階的逐出以 Bélády/MIN 決定，跨階的目的地以成本感知貪婪決定。後者非加權多階快取問題的最佳解，故所得 headroom 為\emph{下界}——真正的最佳解只會更好。此不對稱使 GO 判定安全，而 NO-GO 判定須額外審視。

模擬器的效力取決於兩個前提。其一，成本常數須為實測；程式在讀不到 \S\ref{app:cost} 的量測檔時中止，且模型剖面（GPU 預算、KV 每 token 位元組、成本常數的來源模型）三者必須一致，否則拒絕執行。其二，模擬器須能複現已量測的行為。我們以第 \ref{sec:eval} 節的工作負載比對 \texttt{full\_gpu} 相對 \texttt{cpu\_lru} 的 warm 段成本比：16K 上下文實測 9.00 對模擬 9.73（差 8\%），32K 上下文實測 12.27 對模擬 12.66（差 3\%）。絕對值同樣相符，\texttt{full\_gpu} 為 5{,}033 對 5{,}864\,ms、\texttt{cpu\_lru} 為 559 對 603\,ms。\textbf{驗證僅涵蓋一種工作負載形狀與兩個有鑑別力的上下文長度，本文引用 headroom 的絕對值時一併標明此範圍。}

\paragraph{逐 block 的簡化對前綴結構的工作負載無害。}
vLLM 的查詢語意是連續前綴：\texttt{\_lookup\_complete\_chunks} 回傳一個 token 數，而 token 數無法表達「block 0、1、4、5 命中，2、3 未命中」這種形狀，因此第一個缺口之後的 block 一律重算。我們把模擬器改為此語意並量化其影響：缺口之後的 block 中仍留在任一階的比例為 0.000\% 至 0.464\%（兩條 trace $\times$ 三個策略），其中 61.8\% 至 80.1\% 是首次出現的 block，本來就必須重算。真實流量的重用本身即為前綴結構，Mooncake 的 \texttt{hash\_ids} 就是前綴雜湊，第一個未命中因此恰好落在共用前綴的結尾，缺口之後沒有可供損失的快取內容。此結果同時構成模擬器的一項有效性證據。

% =====================================================================
\bibliographystyle{plain}
\bibliography{refs}

\end{document}
